\documentclass[11pt,oneside]{report}
\usepackage[utf8]{inputenc}
\usepackage[T1]{fontenc}
\usepackage[spanish,es-nodecimaldot]{babel}
\decimalpoint

% Project house style supplied by the author.
% general: author-year bibliography for ML/statistics
% final: no draft overfull-rule markers
% compact: 1.8 cm A4 margins
\usepackage[general,final,compact]{velvetblue}
\addbibresource{references.bib}

% Additional packages not provided by velvetblue.
\usepackage{amssymb,mathtools}
\usepackage{bm}
\usepackage{longtable}
\usepackage{tabularx}
\usepackage{multirow}
\usepackage{siunitx}
\sisetup{detect-all,group-separator={,},group-minimum-digits=4}
\usepackage{float}
\usepackage{placeins}
\usepackage{subcaption}
\usepackage{enumitem}
\usepackage{listings}
\usepackage{fancyhdr}
\usepackage{titlesec}
\usepackage{setspace}
\setstretch{1.08}

% Local report colors complement VelvetBlue without overriding the house style.
\definecolor{GeoGold}{HTML}{D6A128}
\definecolor{GeoGray}{HTML}{5B6573}
\definecolor{GeoLight}{HTML}{F3F6F8}

\hypersetup{
  pdftitle={GeoCebada: Reporte técnico integral},
  pdfauthor={Equipo GeoCebada}
}

\pagestyle{fancy}
\fancyhf{}
\fancyhead[L]{GeoCebada}
\fancyhead[R]{Reporte técnico}
\fancyfoot[C]{\thepage}

\titleformat{\chapter}[display]
{\normalfont\huge\bfseries\color{VelvetBlue}}
{\chaptername\ \thechapter}{12pt}{\Huge}
\titleformat{\section}
{\normalfont\Large\bfseries\color{VelvetBlue}}
{\thesection}{0.7em}{}
\titleformat{\subsection}
{\normalfont\large\bfseries}{\thesubsection}{0.7em}{}

\setlist[itemize]{topsep=0.3em,itemsep=0.2em}
\setlist[enumerate]{topsep=0.3em,itemsep=0.2em}
\lstset{
  basicstyle=\ttfamily\small,
  backgroundcolor=\color{GeoLight},
  frame=single,
  rulecolor=\color{GeoGray},
  breaklines=true,
  columns=fullflexible,
  keepspaces=true,
  showstringspaces=false
}

\newcommand{\artifact}[1]{\nolinkurl{#1}}
\newcommand{\R}{\mathbb{R}}
\newcommand{\Train}{\mathcal{L}}
\newcommand{\Target}{\mathcal{U}}
\newcommand{\Xall}{X_{\Train\cup\Target}}
\newcommand{\ytrain}{y_{\Train}}
\newcommand{\yhidden}{y_{\Target}}
\newcommand{\RMSE}{\operatorname{RMSE}}
\newcommand{\MAE}{\operatorname{MAE}}
\newcommand{\MSE}{\operatorname{MSE}}
\newcommand{\diag}{\operatorname{diag}}

\newcommand{\statusbox}[2]{%
\begin{center}\fcolorbox{VelvetBlue}{GeoLight}{%
\begin{minipage}{0.92\linewidth}\textbf{#1}\quad #2\end{minipage}}
\end{center}}

\newcommand{\safefigure}[4][0.92\linewidth]{%
\begin{figure}[H]\centering
\IfFileExists{#2}{\includegraphics[width=#1]{#2}}{%
\fbox{\parbox{0.86\linewidth}{\centering
Figura no encontrada durante la compilación local.\\
Ruta esperada: \texttt{\detokenize{#2}}}}}
\caption{#3}\label{#4}\end{figure}}

\newenvironment{technicalnote}
{\begin{quote}\small\color{GeoGray}\textbf{Nota técnica.}\ }
{\end{quote}}
\newenvironment{decision}
{\begin{quote}\small\color{VelvetBlue}\textbf{Decisión metodológica.}\ }
{\end{quote}}

\title{\textbf{GeoCebada}\\[0.4em]
Reconstrucción transductiva del rendimiento de cebada\\
\large Reporte técnico integral de ciencia de datos}
\author{Equipo GeoCebada}
\date{\today}
\begin{document}
\pagenumbering{roman}
\maketitle
\begin{abstract}
Este documento consolida el desarrollo técnico de \textbf{GeoCebada} para el Reto
AgroCebada FIRA 2026. El problema operativo consiste en reconstruir
\artifact{RENDIMIENTO_T_HA} para un conjunto fijo de 59 parcelas, disponiendo de 138
parcelas etiquetadas y covariables observables para las 197. La formulación es
\emph{transductiva de objetivo fijo}: los $X$ de las 59 parcelas objetivo pueden intervenir
en transformaciones estrictamente $X$-only, mientras sus rendimientos permanecen ocultos.

Checkpoint~01 fijó la semántica y el contrato reproducible de las fuentes. Checkpoint~02
comprimió series satelitales y fuentes ambientales a una tabla parcelaria de 197 filas y
1,403 predictores. Checkpoint~03 añadió 350 variables agronómicas, 335 variables empíricas
$X$-only, descubrimiento fold-local y benchmarks anidados de representaciones, familias y
ensambles globales.

Checkpoint~04 reformuló la validación para reproducir la información real del reto. La señal
dominante resultó espacial: Moran's $I=0.6797$ con $p=0.001$ y Spearman $\rho=0.4882$ entre
distancia geográfica y diferencia absoluta de rendimiento. La búsqueda local/grafo produjo
\artifact{LocalRidge_C4_all_deterministic_Geo_k24_a30_p1} (Local04D), con RMSE
target-matched medio 0.487845 y pooled 0.495741.

Checkpoint~05 combinó globales y transductivos mediante cross-fitting, stacking y un
mixture-of-experts condicionado en soporte. El mejor mínimo de desarrollo, LocalE13 con
10\% de peso global, alcanzó 0.487243 de RMSE medio y 0.495261 pooled, pero el selector
leave-one-development-split-out obtuvo 0.487934 y 0.495933, ambos ligeramente peores que
Local04D. El gate falló y la confirmación fresca no se puntuó. La fuente canónica posterior
a 05 es \artifact{reports/checkpoint_05/final_predictions.csv} y conserva Local04D.

Después de ese cierre se ejecutó Checkpoint~03D como experimento aditivo de capacidad global.
La configuración final mantuvo los dos protocolos outer de cinco folds y tres folds inner,
pero redujo tuning redundante para evitar un run proyectado de varios días. El benchmark
balanced completó 230 outer fits en 248.085 minutos. Sus finalistas fueron HistGB, XGBoost y
LightGBM, todos sobre G1 agronomic, con RMSE municipality-grouped 0.715339, 0.733080 y
0.748478 respectivamente. El resultado mejora la frontera global grouped, pero no domina
simultáneamente state y grouped ni autoriza comparar directamente esos scores con el RMSE
target-matched de Local04D.

HistGB fue rank científico 1 pero no convertible con el stack ONNX estable. XGBoost rank 2
fue el finalista deployable de mayor rango; el round trip ONNX tuvo diferencia absoluta máxima
$3.814697\times10^{-6}$. Ese ONNX es un estimador global de deployment y no una codificación
exacta del predictor competition-facing Local04D.

La búsqueda global amplia queda cerrada. Cualquier último intento de combinar Local04D con
los nuevos globales 03D debe pasar por un Checkpoint~05B estrecho bajo el mismo protocolo
target-matched; de lo contrario, el trabajo siguiente es integración de submission, app y
deployment.

\end{abstract}
\clearpage
\chapter*{Contrato metodológico global}
\addcontentsline{toc}{chapter}{Contrato metodológico global}
\label{chap:contract}
\statusbox{Regla de oro}{
Las 59 etiquetas ocultas de FIRA nunca se cargan, reconstruyen indirectamente ni se utilizan
para selección. Toda operación que usa $y$ se ajusta exclusivamente con etiquetas visibles
dentro de la partición correspondiente. Las operaciones estrictamente $X$-only pueden usar
las covariables de las 197 parcelas cuando el protocolo transductivo lo permite.
}

\section*{Universo estadístico}
Sea $\Train$ el conjunto de 138 parcelas con rendimiento observado y $\Target$ el conjunto
de 59 parcelas de predicción. Para cada parcela $i$ se observa $x_i\in\R^p$ y, si
$i\in\Train$, un rendimiento $y_i$. El objeto de inferencia es
\[
\widehat{\bm y}_{\Target}=(\hat y_i:i\in\Target),
\]
comparado por el evaluador con el vector reservado $\bm y_{\Target}$.

La información disponible durante el desarrollo es
\[
\mathcal I=
\{X_{\Train},\bm y_{\Train},X_{\Target},G,E_{\mathrm{public}}\},
\]
donde $G$ denota geometría/estructura espacial y $E_{\mathrm{public}}$ denota evidencia
externa pública permitida. Por construcción,
\[
\bm y_{\Target}\notin\mathcal I.
\]

\section*{Inducción, transducción y por qué la distinción importa}
En un problema inductivo convencional se pretende aprender una función
$f:\mathcal X\rightarrow\mathbb R$ que generalice a puntos futuros no conocidos. Aquí el
conjunto de consultas $X_{\Target}$ está fijado de antemano. Esto habilita transformaciones
transductivas dependientes de la geometría conjunta de $\Xall$: estandarización $X$-only,
PCA $X$-only, densidad, vecindarios, grafos y métricas de soporte. La restricción es
\[
T_X=T(X_{\Train},X_{\Target})\quad\text{es admisible si no usa }\bm y_{\Target},
\]
mientras cualquier transformación supervisada debe recalcularse usando exclusivamente las
etiquetas permitidas por el fold.

\section*{Métricas}
Para un conjunto de evaluación $S$,
\begin{align}
\RMSE(S)&=\sqrt{\frac{1}{|S|}\sum_{i\in S}(y_i-\hat y_i)^2},\\
\MAE(S)&=\frac{1}{|S|}\sum_{i\in S}|y_i-\hat y_i|,\\
R^2(S)&=
1-\frac{\sum_{i\in S}(y_i-\hat y_i)^2}
{\sum_{i\in S}(y_i-\bar y_S)^2}.
\end{align}

La selección principal del Checkpoint~04 usa RMSE medio entre pseudo-competiciones
target-matched. También se conservan RMSE pooled, MAE, peor split y protocolos de estrés,
porque con $n=138$ las conclusiones pueden cambiar si un split espacial domina el promedio.

\section*{Jerarquía de evidencia}
La interpretación final sigue esta prioridad:
\begin{enumerate}
\item invariantes y semántica de fuentes;
\item predicciones pseudo-target realmente out-of-sample respecto a $y$;
\item estabilidad entre splits y protocolos;
\item robustez espacial/grupada;
\item diagnósticos de soporte, discrepancia y plausibilidad;
\item sólo al final, predicciones de las 59 parcelas, sin utilizarlas para ajustar reglas.
\end{enumerate}
Esta jerarquía evita seleccionar modelos por el aspecto visual de sus 59 valores finales,
optimizar hiperparámetros contra un único protocolo conveniente o reinterpretar información
pública municipal como si fuese etiqueta parcelaria.

\clearpage
\tableofcontents
\clearpage
\listoftables
\clearpage
\listoffigures
\clearpage
\pagenumbering{arabic}
\chapter{Checkpoint 01: Fundación de datos, semántica y contrato reproducible}
\label{chap:checkpoint01}
\section{01A. Unidad de análisis, objetivo y frontera de información}
\label{sec:checkpoint01a}
\subsection{Contrato del reto y unidad estadística}

La primera decisión técnica fue separar con precisión la \emph{unidad de observación de las
fuentes} de la \emph{unidad supervisada de predicción}. La unidad supervisada es la parcela,
identificada de forma única por \artifact{ID_POLIGONO}. El archivo oficial de
rendimiento/split contiene 197 parcelas: 138 filas de \artifact{ENTRENAMIENTO} con
\artifact{RENDIMIENTO_T_HA} observado y 59 filas de \artifact{PREDICCION} con el objetivo
reservado por FIRA.

\begin{table}[H]
\centering
\caption{Contrato básico del problema supervisado.}
\label{tab:cp01-contract}
\begin{tabular}{lr}
\toprule
Elemento & Valor \\
\midrule
Parcelas totales & 197 \\
Parcelas etiquetadas & 138 \\
Parcelas objetivo & 59 \\
Estados & Hidalgo, Puebla, Tlaxcala \\
Identificador canónico & \artifact{ID_POLIGONO} \\
Variable objetivo & \artifact{RENDIMIENTO_T_HA} \\
Ciclo objetivo & abril--octubre de 2025 \\
\bottomrule
\end{tabular}
\end{table}

El hecho de que BASIC y PRO contengan decenas de miles de filas no aumenta el tamaño muestral
supervisado. Una fila BASIC corresponde conceptualmente a una observación
parcela--fecha--sensor; una fila PRO tiene la misma naturaleza longitudinal. Por tanto,
\[
n_{\mathrm{supervisado}}=138,
\]
no $107{,}666$ ni $47{,}804$. Tratar las filas temporales como observaciones independientes
con la etiqueta parcelaria replicada produciría pseudorreplicación severa, correlación
intra-parcela ignorada y leakage entre train y validación. Toda la arquitectura de datos
posterior conserva esta distinción.

\subsection{Objetivo estadístico inicial y evolución posterior}

Durante los Checkpoints~01--03 se mantuvo una formulación convencional de regresión:
construir una representación parcelaria $x_i$ y estudiar modelos
\[
\hat y_i=f_\theta(x_i)
\]
bajo validación espacialmente consciente. Esta etapa era necesaria para establecer qué
fuentes, representaciones y familias de modelos eran razonables antes de explotar la
estructura fija del conjunto objetivo.

Posteriormente, en el Checkpoint~04, la doctrina se refinó a una formulación
transductiva. Esa reformulación no invalida los checkpoints previos: los convierte en
evidencia de ingeniería, representación y baseline. La diferencia central es que en el
problema final conocemos de antemano los 59 vectores $X_{\Target}$ que serán evaluados. Por
ello, el objeto final no es una función universal para parcelas futuras sino el vector
\[
\widehat{\bm y}_{\Target}\in\R^{59}.
\]

\subsection{Frontera de información}

Desde el primer checkpoint se mantuvo una separación estricta entre:
\begin{enumerate}
\item información observable y permitida;
\item información derivada de etiquetas visibles;
\item etiquetas reservadas.
\end{enumerate}

Para los 59 objetivos:
\[
X_{\Target}\;\text{es observable},\qquad
\bm y_{\Target}\;\text{es no observable}.
\]
La regla se aplica no sólo al fit final sino también a selección de variables,
hiperparámetros, elección de vecinos, thresholds de routing y decisiones de ensamble.

En los checkpoints convencionales, transformaciones entrenables como imputación, escalado,
PCA, PLS y descubrimiento supervisado se ajustan dentro de cada fold. En la fase transductiva
se permite que transformaciones estrictamente $X$-only conozcan $X_{\Target}$, pero nunca
$\bm y_{\Target}$. Esta distinción se vuelve especialmente importante para grafos,
estandarización transductiva, densidad y selección de pseudo-competiciones.

\subsection{Horizonte de información: clean frente a competition}

La ingeniería de datos reconoció desde el principio que una variable puede ser válida para la
competencia y, sin embargo, no corresponder a un sistema de pronóstico prospectivo. Por eso se
crearon dos modos conceptuales.

\paragraph{clean.}
Incluye información histórica, estática o disponible bajo una interpretación prospectiva
conservadora. Su objetivo es permitir una lectura científica de generalización sin usar
resultados contemporáneos que sólo aparecen después del ciclo.

\paragraph{competition.}
Incluye el conjunto clean más información pública/transductiva permitida para reconstruir el
rendimiento de las 59 parcelas después de observar el ciclo 2025. Ejemplos: series 2025
completas, WaPOR 2025, clima 2025 y, posteriormente, SIAP 2025 como evidencia municipal
contemporánea.

La separación no equivale a \emph{seguro} frente a \emph{leakage}. Una variable competition es
válida si procede de una fuente pública permitida y no codifica los 59 rendimientos ocultos.
La categoría existe para conservar honestidad semántica sobre el momento en el que la
información estaría disponible en un escenario de despliegue distinto.

\subsection{Implicación small-\texorpdfstring{$n$}{n}/high-\texorpdfstring{$p$}{p}}

Desde el inicio se anticipó que la tabla final sería de dimensión elevada respecto al número
de etiquetas. Esta condición condiciona todo el proyecto:
\[
n_{\Train}=138,\qquad p\gg n.
\]
Consecuencias:
\begin{itemize}
\item cualquier selección guiada por $y$ tiene alta varianza;
\item un grid de hiperparámetros muy grande puede sobreoptimizar los folds;
\item los errores de validación deben leerse junto con dispersión y protocolos espaciales;
\item Ridge y PLS son baselines fundamentales, no modelos demasiado simples;
\item árboles/boosting requieren fuerte disciplina de validación;
\item una mejora de milésimas en el mismo conjunto de desarrollo no debe interpretarse como
ganancia real sin un control split-excluded.
\end{itemize}

Esta última regla resultó decisiva en Checkpoint~04F: un blend Local/Graph obtuvo un mínimo
ligeramente inferior al LocalRidge en la tabla completa de desarrollo, pero no se promovió
porque la selección leave-one-split-out no sostuvo la ventaja.

\subsection{Definición operacional de éxito}

El proyecto no busca maximizar una métrica abstracta de representación. Cada checkpoint debe
reducir incertidumbre sobre una decisión concreta. En las primeras fases, la pregunta es si
una fuente puede integrarse sin ambigüedad y si una representación mantiene señal bajo
validación. En la fase transductiva, el criterio se hace más estricto:

\begin{quote}
¿La información observable de las 197 parcelas permite reducir el error en pseudo-competiciones
que reproducen la geometría de las 59 parcelas reales, sin acceder en ningún momento a sus
etiquetas ocultas?
\end{quote}

Esta formulación conecta ingeniería de datos, estadística espacial, aprendizaje supervisado y
semi-supervisado en un único contrato de decisión.

\section{01B. Inventario multifuente y semántica de los datos}
\label{sec:checkpoint01b}
\subsection{Fuentes oficiales}

El sistema integra sensores remotos, geometría, clima, topografía y estadísticas públicas.
La \cref{tab:cp01-sources} resume el inventario principal congelado en el contrato de
datos.

\begin{longtable}{p{0.20\linewidth}p{0.19\linewidth}p{0.49\linewidth}}
\caption{Fuentes principales y semántica operacional.}
\label{tab:cp01-sources}\\
\toprule
Fuente & Cobertura & Uso principal \\
\midrule
\endfirsthead
\toprule
Fuente & Cobertura & Uso principal \\
\midrule
\endhead
Split oficial & 197 parcelas & ID, área, conjunto y rendimiento observado para 138 parcelas.\\
Geometrías oficiales & 197 polígonos, EPSG:4326 & centroides, distancias, extracción espacial y QA administrativa.\\
BASIC & $107{,}666\times96$; 2022--2025 & Sentinel-2/Landsat, fenología histórica y 2025, índices de vegetación/agua/estrés.\\
PRO & $47{,}804\times20$; 2025 & Planet NDVI/EVI/LAI/MSAVI y estadísticos temporales.\\
Clima oficial & 48 meses por variable, 2022--2025 & precipitación, Tmin, Tmax, climatología/anomalías.\\
Topografía oficial & 120 m & elevación y pendiente.\\
INEGI municipios & Hidalgo/Puebla/Tlaxcala & CVEGEO, unión SIAP y QA espacial.\\
CHIRPS v3 & 214 días, abr--oct 2025 & precipitación diaria y episodios secos/extremos; excluido en varias fases por QC.\\
WaPOR v3 & 21 dekads, abr--oct 2025 & AETI y NPP; balance hídrico/productividad.\\
SoilGrids & 9 propiedades $\times$ 4 profundidades & perfil edáfico estático y heterogeneidad vertical.\\
SIAP & 2003--2025 & contexto municipal histórico y proxy contemporáneo 2025.\\
INEGI CEM & $\sim15$ m & elevación/relieve de mayor resolución.\\
\bottomrule
\end{longtable}

\subsection{BASIC: representación longitudinal multiespectral}

BASIC contiene 107,666 registros y 96 columnas. Los registros cubren las 197 parcelas entre
2022-01-01 y 2025-12-31 y provienen de Landsat y Sentinel-2. Las fechas se interpretan
estrictamente como \artifact{DD/MM/YYYY}; la llave parcela--fecha--sensor es única.

La estructura de variables comprende 23 familias de índices con cuatro resúmenes espaciales
por observación ---promedio, desviación estándar, máximo y mínimo---, además de metadatos de
fecha, sensor y nubosidad. Dos decisiones son esenciales:
\begin{enumerate}
\item no se mezclan silenciosamente variables homónimas entre sensores;
\item la ausencia sistemática de ciertos índices por sensor no se corrige con un fill global,
porque esa ausencia es estructural y contiene provenance.
\end{enumerate}

En particular, fórmulas provider-specific como MRA y VI6T no se inventan. Se conserva su
semántica tal como aparece en el suministro original. Esta política es importante porque un
índice de nombre opaco no debe recibir una interpretación espectral creada a posteriori sólo
para justificar una feature.

\subsection{PRO: señal Planet 2025}

PRO aporta 47,804 registros por 20 columnas, todas las parcelas y fechas de 2025-01-01 a
2025-12-30. El sensor Planet contiene NDVI, EVI, LAI y MSAVI con estadísticos espaciales.
Esta fuente es particularmente relevante para capturar el ciclo contemporáneo con mayor
densidad/resolución, pero sigue siendo una colección de observaciones longitudinales. Su
incorporación supervisada ocurre sólo después de agregarla de forma parcelaria o de construir
métricas temporales que respetan la entidad parcela.

La coexistencia de Sentinel-2 y Planet permite además construir features de concordancia
inter-sensor. Esa comparación no presupone equivalencia radiométrica: se utiliza para
cuantificar si ambas cadenas describen una forma temporal compatible para la misma parcela.

\subsection{Clima oficial, CHIRPS y WaPOR}

La fuente climática oficial contiene 48 rasters mensuales para cada una de precipitación,
temperatura mínima y temperatura máxima, con continuidad 2022--2025 y CRS EPSG:4326. A partir
de ella se construyen climatologías históricas, totales estacionales y anomalías 2025.

CHIRPS v3 aporta precipitación diaria para 214 fechas entre abril y octubre de 2025. Aunque su
resolución temporal permite construir número de días de lluvia, máximos móviles y duración de
episodios secos, el pipeline mantuvo CHIRPS fuera de varias representaciones de modelado porque
Feature Table v1 sólo retuvo una característica y se decidió no sobreinterpretar una
extracción que requería QC adicional.

WaPOR v3 aporta AETI y NPP en 21 periodos dekadales. Dado que sus bandas representan tasas,
una agregación estacional correcta multiplica por la duración del dekad:
\begin{equation}
AETI_{\mathrm{season}}
=
\sum_{k=1}^{21} AETI_k d_k.
\label{eq:wapor-duration}
\end{equation}
La misma regla se aplica a otros productos de tasa. Esta distinción impide sumar valores con
unidades por día como si fueran acumulados. La capa agronómica utilizará posteriormente estas
magnitudes para construir proxies de productividad hídrica \citep{Dhonthi2024}.

\subsection{Suelo y topografía}

SoilGrids aporta 36 GeoTIFF: nueve propiedades por cuatro profundidades:
pH, arcilla, arena, limo, carbono orgánico, nitrógeno, capacidad de intercambio catiónico,
densidad aparente y fracción de fragmentos gruesos. Los TIFF descargados mediante WCS pueden
carecer de CRS en el header; sólo los productos positivamente identificados reciben el
fallback ESRI:54052 definido por el contrato. Las conversiones de escala INT16 a unidades
físicas se aplican según la especificación del producto, nunca por inferencia ad hoc.

La topografía combina el producto oficial de 120 m y el CEM estatal de INEGI de
aproximadamente 15 m. Las fuentes se mantienen separadas mediante prefijos para evitar que una
resolución sustituya silenciosamente a otra. En modelos locales, la elevación y la pendiente
pueden actuar como contexto ambiental; sin embargo, la distancia geográfica directa terminó
teniendo una asociación más fuerte con diferencias de rendimiento que la distancia
multivariada completa.

\subsection{SIAP y la llave administrativa}

El archivo histórico SIAP cubre 2003--2025. La llave municipal estable es
\begin{equation}
\mathrm{CVEGEO}
=
\mathrm{zfill}_2(\mathrm{Idestado})
\Vert
\mathrm{zfill}_3(\mathrm{Idmunicipio}),
\label{eq:cvegeo}
\end{equation}
donde $\Vert$ denota concatenación. Nunca se une únicamente por
\artifact{Idmunicipio}, porque dicho identificador no es globalmente único entre estados.

Se identificaron al menos las categorías \artifact{Cebada grano} y
\artifact{Cebada forrajera en verde}, además de ciclos Otoño--Invierno y
Primavera--Verano y modalidades Riego/Temporal. Mezclar grano y forraje modificaría la
semántica de la respuesta, por lo que se prohíbe explícitamente.

En Feature Table v1 se preservó una versión histórica del agregador SIAP. En
Checkpoint~04E.1 se volvió a la tabla detallada para filtrar de forma explícita
\artifact{Cebada grano + Primavera-Verano + Temporal + CVEGEO + 2025}.
La distinción entre el agregado legado y el scope exacto es importante para interpretar por
qué una fuente administrativamente correcta puede seguir siendo un proxy pobre del rendimiento
parcelario. Debe señalarse además que la elección de Primavera--Verano se alineó con el ciclo
abril--octubre del reto; no se presenta aquí como una certificación externa independiente de
semántica de ciclo.

\subsection{Asignación municipal y casos de borde}

La etiqueta administrativa oficial de parcela se preserva como referencia para joins. La
superposición con polígonos INEGI se utiliza como QA/fallback y no sustituye silenciosamente
el dato oficial. Se documentaron, entre otros, AGC\_020 con una pequeña intersección
interestatal, AGC\_129 atravesando municipios en Hidalgo y AGC\_048 con discrepancia entre
etiqueta oficial y overlay. Esta política evita que un algoritmo de overlay convierta
automáticamente un caso geométrico limítrofe en una reasignación administrativa.

\subsection{Fuentes investigadas pero no promovidas}

ERA5-Land fue explorado y se implementó tooling de descarga, pero las solicitudes a CDS
presentaron problemas repetidos de costo/backend y no se obtuvo un artefacto completo
certificado para el pipeline. En lugar de permitir una fuente parcial, se declaró opcional y
se excluyó de downstream. Esta decisión ilustra una regla general: una fuente potencialmente
útil no entra al modelo por prestigio o disponibilidad nominal; debe tener cobertura,
unidades y proceso de extracción auditables.

\section{01C. Auditoría, contrato de datos e integración reproducible}
\label{sec:checkpoint01c}
\subsection{Data Contract v2}

La etapa de datos se cerró mediante un contrato ejecutable, no mediante una inspección manual
única. Los artefactos centrales son \artifact{configs/data_contract_v2.yaml} y
\artifact{tools/audit_data_contract_v2.py}. El audit completo produjo:
\[
46\ \text{PASS},\qquad
0\ \text{WARN},\qquad
0\ \text{FAIL},\qquad
1\ \text{SKIP}.
\]

Entre las invariantes comprobadas están:
\begin{itemize}
\item 197 IDs oficiales únicos;
\item partición exacta 138/59;
\item target visible sólo en entrenamiento;
\item formas BASIC $107{,}666\times96$ y PRO $47{,}804\times20$;
\item cobertura de las 197 parcelas en ambas tablas satelitales;
\item parsing estricto de fechas y unicidad parcela--fecha--sensor;
\item validez geométrica y CRS;
\item continuidad de rasters climáticos;
\item continuidad diaria CHIRPS;
\item metadatos/unidades WaPOR;
\item inventario SoilGrids y reglas de CRS;
\item cobertura SIAP 2003--2025, schemas y duplicados;
\item inventario CEM;
\item exclusión explícita de ERA5-Land como fuente no certificada.
\end{itemize}

El contrato funciona como \emph{gate}: si cambian fuentes crudas, primero debe volver a
ejecutarse el audit; ningún checkpoint posterior debería asumir que una descarga nueva conserva
las mismas unidades o cobertura.

\subsection{Inmutabilidad de fuentes y transformación versionada}

El repositorio separa:
\begin{enumerate}
\item evidencia fuente inmutable bajo \artifact{data/source/};
\item datos externos reproducibles/locales bajo \artifact{data/raw/external/};
\item artefactos derivados pequeños y versionados;
\item reportes generados por checkpoints.
\end{enumerate}

Las correcciones no se escriben sobre el archivo fuente. Por ejemplo, metadata defectuosa de
un NetCDF puede normalizarse en un artefacto derivado, pero el input original permanece
intacto. Esto permite reconstruir la cadena
\[
\text{source}
\longrightarrow
\text{loader}
\longrightarrow
\text{normalización}
\longrightarrow
\text{feature}
\longrightarrow
\text{modelo}.
\]

La separación tiene una consecuencia estadística: cualquier resultado puede auditarse contra
la versión exacta de la tabla que lo generó, reduciendo el riesgo de que un mismo nombre de
feature represente transformaciones distintas en checkpoints diferentes.

\subsection{Integration fixtures}

Se creó un fixture real de 12 parcelas con IDs compartidos entre fuentes. La selección incluye
dos parcelas por estrato Estado$\times$CONJUNTO, diversidad target-free y casos de borde
administrativos. El fixture contiene
\[
12\ \text{parcelas},\qquad
192\ \text{filas BASIC},\qquad
96\ \text{filas PRO},\qquad
201\ \text{filas SIAP},
\]
además de geometrías, clima, CHIRPS, topografía, SoilGrids, WaPOR y CEM. El conjunto ocupa
aproximadamente 3.77 MiB y permite probar joins/extractores sin requerir el archivo externo
completo.

La selección del fixture no usa \artifact{RENDIMIENTO_T_HA}. Esto es relevante porque incluso
un dataset de integración podría introducir un sesgo sutil si se eligieran ejemplos por
valores extremos de $y$.

\subsection{CRS y operaciones espaciales}

Las operaciones espaciales se realizan con una política explícita de CRS. EPSG:4326 es
adecuado para almacenamiento/intercambio de coordenadas angulares, pero no para interpretar
distancias euclidianas en grados como kilómetros. Por ello, cualquier cálculo métrico debe
transformarse o usar una fórmula geodésica apropiada.

En etapas posteriores, la componente geográfica de las distancias transductivas se deriva de
centroides parcelarios y se normaliza antes de combinarse con distancias agronómicas. La
consistencia de CRS es especialmente importante porque la geografía terminó siendo la señal
de transferencia más fuerte del proyecto.

\subsection{Missingness como semántica y no sólo como problema numérico}

La política de datos distingue tres casos:
\begin{enumerate}
\item valor ausente porque la fuente no cubre una parcela/fecha;
\item variable no definida para un sensor;
\item valor eliminado por QA.
\end{enumerate}
No todos los NaN significan lo mismo. Durante la construcción se eliminan columnas
completamente nulas porque no pueden aportar información. Las columnas parcialmente faltantes
se conservan y la imputación aprendida se delega al pipeline de modelado. De esta forma, una
decisión de ingeniería no utiliza accidentalmente distribución de validación.

\subsection{Hashing, manifest y trazabilidad}

Los builds derivados registran hashes de inputs, commit de Git y manifest de features. Esta
trazabilidad permite responder tres preguntas que son particularmente importantes en una
competencia con múltiples fuentes:
\begin{enumerate}
\item ¿qué archivo crudo originó una variable?;
\item ¿qué versión de código produjo la tabla usada por un modelo?;
\item ¿qué features pertenecían realmente al conjunto clean o competition en ese commit?
\end{enumerate}

La respuesta no depende de memoria humana ni de nombres de notebooks. El estado reproducible
queda materializado en YAML, JSON y CSV versionados.

\subsection{Criterios de cierre del Checkpoint 01}

Checkpoint~01 no produjo todavía un modelo ni una tabla final $197\times p$. Su cierre
significó algo más fundamental: cada fuente relevante tenía una semántica, una política de
unidades, una llave de unión, un tratamiento de missingness y un mecanismo de auditoría
documentados. El Checkpoint~02 pudo entonces operar sobre una base reproducible en vez de
incorporar limpieza implícita dentro del código de modelado.

\begin{decision}
Una feature sólo puede entrar a modelado si puede trazarse hasta una fuente y una regla de
transformación. La ausencia de provenance se trata como un defecto de datos, no como un
detalle de documentación.
\end{decision}

\chapter{Checkpoint 02: Tabla parcelaria, diagnóstico de covariables y validación congelada}
\label{chap:checkpoint02}
\section{02A. Feature Table v1 y separación clean/competition}
\label{sec:checkpoint02a}
\subsection{Objetivo de la compresión parcelaria}

Feature Table v1 transforma fuentes heterogéneas y longitudinales en una representación
determinista de una fila por parcela. El artefacto no es un dataset limpio en el sentido
de haber aplicado imputación global o escalado aprendido; es una tabla \emph{model-ready}
porque resuelve las uniones, unidades, agregaciones y provenance, dejando cualquier
transformación entrenable dentro de la validación.

El build validado produjo:
\begin{table}[H]
\centering
\caption{Dimensiones de Feature Table v1.}
\label{tab:cp02-featuretable}
\begin{tabular}{lrr}
\toprule
Artefacto & Filas & Columnas \\
\midrule
\artifact{parcel_features_clean.csv} & 197 & 699 \\
\artifact{parcel_features_competition.csv} & 197 & 1,408 \\
\artifact{parcel_features_all.csv} & 197 & 1,408 \\
\bottomrule
\end{tabular}
\end{table}

El manifest contiene 1,403 predictores: 694 clean y 709 competition-only. Las columnas
adicionales corresponden a ID/metadata/target/split. Todas las tablas preservan exactamente
197 IDs, 138 targets observados y 59 targets missing.

\subsection{Familias de variables}

El inventario de Checkpoint~02 clasifica las features por source, family y mode. Entre las
features clean aparecen 539 BASIC históricos, 90 SoilGrids, 24 de clima oficial histórico,
14 SIAP históricos, ocho de topografía, seis geométricas y cinco CEM, además de campos
administrativos. La capa competition añade 539 BASIC 2025, 97 PRO, 45 clima 2025, 22 WaPOR,
cinco SIAP contemporáneas y la feature CHIRPS disponible.

Esta taxonomía permite realizar ablaciones semánticas. No basta con preguntar si una columna
mejora RMSE; interesa identificar si el incremento proviene de geometría, clima, historial
satelital, señal contemporánea o una estadística administrativa.

\subsection{Agregación temporal}

Para una serie longitudinal $x_{i,s,t}$ de parcela $i$, sensor $s$ y fecha $t$, Feature Table
v1 genera resúmenes compatibles con el source contract: históricos, 2025, mensuales,
estacionales y por sensor. La operación conceptual es
\begin{equation}
x_{i,s,\mathcal T}^{(g)}
=
g\left(\{x_{i,s,t}:t\in\mathcal T\}\right),
\label{eq:feature-aggregation}
\end{equation}
donde $g$ puede ser media, desviación, mínimo, máximo u otro resumen explícito.

La agregación no cruza sensores por conveniencia. Por ejemplo, Sentinel-2 NDVI y Landsat NDVI
son variables distintas aun si comparten nombre del índice. Esta decisión evita que
disponibilidad temporal/sensor-specific missingness se convierta en un promedio sin
interpretación física.

La tabla final no pretende agotar el contenido temporal. El Checkpoint~04A vuelve a
trayectorias mensuales para construir correlación, lag y DTW. En otras palabras, Feature Table
v1 es una proyección tabular de los datos longitudinales, no una afirmación de suficiencia
estadística.

\subsection{Invariantes numéricas}

El build elimina features completamente nulas y verifica ausencia de infinitos. Esto es
diferente de imputar missingness de forma supervisada. Sea $X$ la matriz derivada. La etapa de
datos puede retirar una columna $j$ si
\[
\forall i,\; X_{ij}=\mathrm{NaN},
\]
porque no contiene información en ninguna parcela. En cambio, para una columna parcialmente
observada, la imputación debe estimarse en el conjunto de entrenamiento del fold:
\[
\tilde X_{ij}=
\begin{cases}
X_{ij}, & X_{ij}\text{ observado},\\
\widehat m_j^{(\mathrm{train})}, & X_{ij}\text{ missing}.
\end{cases}
\]
Usar una mediana calculada con el fold de validación no usa $y$, pero en una evaluación
inductiva puede transferir información de distribución. Por disciplina, Checkpoints~02--03
mantuvieron el preprocessing aprendido dentro del pipeline.

\subsection{Representación y provenance}

Cada entrada del manifest registra nombre de columna, fuente, familia y modo. El manifest
funciona como un esquema semántico que permite formar subconjuntos sin depender de regex
frágiles sobre miles de nombres.

Esta estructura fue crítica en Checkpoint~03, donde se construyeron capas agronómicas y
empíricas separadas, y en Checkpoint~04, donde C0, C1 y C4 se resolvieron mediante provenance
en vez de mantener listas manuales de columnas.

\subsection{Por qué no escalar el target}

En discusiones tempranas se consideró multiplicar el rendimiento por 100 para evitar RMSE
numéricamente menor que uno. Esa transformación no mejora el problema. Si
\[
y_i'=c\,y_i,\qquad \hat y_i'=c\,\hat y_i,
\]
entonces
\[
\RMSE(y',\hat y')=|c|\,\RMSE(y,\hat y).
\]
La escala cambia la magnitud reportada, pero no el orden relativo de modelos ni la calidad de
la reconstrucción. Por ello se conserva la unidad natural de toneladas por hectárea.

\subsection{Condición de congelamiento}

Feature Table v1 queda congelada como artefacto reproducible. Las capas posteriores son
\emph{aditivas}: no reescriben las columnas base ni modifican retrospectivamente un build para
hacer que un resultado posterior parezca parte de la representación original. Esta política
hace posible comparar B0, B1, B2, etc. bajo la misma fuente de verdad.

\section{02B. Diagnóstico train--prediction, soporte y ablaciones}
\label{sec:checkpoint02b}
\subsection{Train versus prediction}

Una vez congelada la tabla, el primer análisis no supervisado comparó la distribución de las
138 filas etiquetadas con las 59 filas objetivo. Se examinaron 1,401 variables numéricas. Para
cada feature se calcularon, entre otros, diferencia estandarizada de medias, missingness y
fracción del conjunto objetivo fuera del rango observado de entrenamiento.

Una forma genérica del standardized mean difference es
\begin{equation}
\mathrm{SMD}_j
=
\frac{\bar x_{j,\Train}-\bar x_{j,\Target}}
{\sqrt{\frac{s^2_{j,\Train}+s^2_{j,\Target}}{2}}}.
\label{eq:smd}
\end{equation}
Sólo seis variables superaron al menos un threshold descriptivo. Cuatro correspondían a
resúmenes históricos Sentinel-2 de agosto, una a la desviación de densidad aparente
SoilGrids 30--60 cm y una a EVI mínimo Sentinel-2 2025.

\begin{table}[H]
\centering
\caption{Principales flags descriptivos del diagnóstico train--prediction.}
\label{tab:cp02-shift}
\small
\begin{tabular}{p{0.60\linewidth}rr}
\toprule
Feature & $|\mathrm{SMD}|$ & fuera de rango \\
\midrule
Sentinel-2 ENDDI histórico, agosto & 0.5343 & 0.0000 \\
Sentinel-2 DSWI2 histórico, agosto & 0.5303 & 0.0000 \\
Sentinel-2 MCRC histórico, agosto & 0.5296 & 0.0169 \\
Sentinel-2 CRC histórico, agosto & 0.5035 & 0.0339 \\
SoilGrids bdod 30--60 cm, std & 0.3218 & 0.1017 \\
Sentinel-2 EVI mínimo 2025, std & 0.2212 & 0.1017 \\
\bottomrule
\end{tabular}
\end{table}

El resultado sugería ausencia de un shift univariado masivo, pero no descartaba shift
multivariado. Esa pregunta fue retomada formalmente con clasificación adversarial en
Checkpoint~04A.

\subsection{Ablaciones estructuradas}

En lugar de entrenar inmediatamente sobre las 1,401 variables numéricas, Checkpoint~02 definió
una secuencia acumulativa de ocho representaciones:

\begin{table}[H]
\centering
\caption{Ablaciones de Checkpoint 02.}
\label{tab:cp02-ablations}
\small
\begin{tabular}{llr}
\toprule
ID & Contenido & $p$ \\
\midrule
A0 & geometría + administración & 12 \\
A1 & A0 + ambiente estático & 115 \\
A2 & A1 + historia SIAP & 129 \\
A3 & A2 + clima histórico & 153 \\
A4 & A3 + BASIC histórico & 692 \\
A5 & clean completo & 692 \\
A6 & clean + remote/weather 2025, sin SIAP outcome & 1,396 \\
A7 & competition completo, incluido SIAP 2025 & 1,401 \\
\bottomrule
\end{tabular}
\end{table}

A4 y A5 coinciden numéricamente porque, bajo el manifest de esa versión, todas las features
clean numéricas ya estaban contenidas en A4. La duplicación se mantuvo como checkpoint
semántico: A4 representa hasta BASIC histórico y A5 todo clean.

\subsection{Por qué las ablaciones importan en high-dimensional small-\texorpdfstring{$n$}{n}}

Cuando $p\gg n$, agregar una familia completa puede degradar un modelo incluso si contiene
señal real. La varianza del ajuste, colinealidad y ruido pueden dominar. Por ello se mide
\begin{equation}
\Delta\RMSE_{A\rightarrow B}
=
\RMSE(B)-\RMSE(A)
\label{eq:ablation-delta}
\end{equation}
por modelo y protocolo, en vez de suponer monotonicidad con el número de features.

El ejemplo más claro aparece con Ridge: en state-stratified, A3 obtiene RMSE medio 0.5298,
mientras A4/A5 suben a 0.6707. Agregar cientos de resúmenes satelitales a un modelo lineal de
alpha fijo no es una mejora automática. ExtraTrees es mucho menos sensible en ese protocolo:
A3 da 0.5035 y A7 0.4966.

\subsection{Clean y competition como ejes experimentales}

Checkpoint~02 mostró que la señal 2025 completa podía ayudar en algunos modelos, pero también
que SIAP contemporáneo no tenía por qué mejorar de forma automática. ExtraTrees pasa de
0.5001 en A6 a 0.4966 en A7 bajo state-stratified, una ganancia pequeña; en
municipality-grouped, A6=0.8835 y A7=0.8916, es decir, el agregado completo empeora. Este
patrón anticipa la conclusión mucho más fuerte de Checkpoint~04E.1: una fuente contemporánea
puede ser legal y plausible sin ser un buen proxy parcelario.

\subsection{Rango y extrapolación}

La fracción de objetivos fuera del rango train,
\[
q_j=
\frac{1}{|\Target|}
\sum_{i\in\Target}
\mathbf 1\{x_{ij}<\min_{\Train}x_{\cdot j}\;\vee\;
x_{ij}>\max_{\Train}x_{\cdot j}\},
\]
es útil como alerta, pero no prueba que una parcela sea imposible de predecir. Un punto puede
estar fuera del rango de una feature marginal y, simultáneamente, cerca de entrenamiento en
una representación multivariada relevante. Por esa razón el soporte final de Checkpoint~04
combina geografía, embeddings, temporalidad y clasificación adversarial.

\begin{technicalnote}
Los flags de covariate shift de Checkpoint~02 son descriptivos. No se usaron para inferir
pseudo-etiquetas ni para eliminar variables usando conocimiento de las 59 respuestas.
\end{technicalnote}

\section{02C. Protocolos de validación y modelos base}
\label{sec:checkpoint02c}
\subsection{Dos protocolos de validación}

Checkpoint~02 congeló cinco folds para dos objetivos distintos.

\paragraph{fold state stratified.}
Preserva aproximadamente la composición por estado en cada fold. Es un protocolo de
interpolación más cercano a parcelas nuevas dentro de regiones representadas.

\paragraph{fold municipality grouped.}
Mantiene juntas las parcelas de un mismo municipio, por lo que un municipio retenido no
aparece parcialmente en train y validación. Este protocolo es deliberadamente más severo y
mide sensibilidad espacial/administrativa.

La diferencia entre ambos no se considera un defecto del benchmark. Es evidencia sobre el
tipo de generalización que cada modelo logra. Si
\[
\RMSE_{\mathrm{grouped}}\gg\RMSE_{\mathrm{state}},
\]
la explicación natural es que una parte de la capacidad predictiva depende de contexto local
que desaparece al retener municipios completos.

\subsection{Baselines deliberadamente simples}

La fase inicial evaluó DummyMean, Ridge de alpha fijo y ExtraTrees regularizado. La imputación
por mediana se ajustó dentro de cada fold. No se hizo PCA global, selección supervisada global,
Optuna ni CatBoost en esta etapa.

Ridge resuelve
\begin{equation}
\widehat{\bm\beta}
=
\arg\min_{\bm\beta}
\left\{
\sum_{i\in\Train_f}
(y_i-\beta_0-x_i^\top\bm\beta)^2
+\lambda\|\bm\beta\|_2^2
\right\},
\label{eq:ridge-global}
\end{equation}
lo que ofrece una referencia lineal regularizada en presencia de colinealidad
\citep{HoerlKennard1970}. ExtraTrees proporciona un contraste no lineal basado en árboles
altamente aleatorizados \citep{Geurts2006}.

\subsection{Resultados seleccionados}

\begin{table}[H]
\centering
\caption{Resultados representativos del Checkpoint 02.}
\label{tab:cp02-baselines}
\small
\begin{tabular}{llrr}
\toprule
Protocolo & Modelo/ablación & RMSE medio & $R^2$ medio \\
\midrule
state-stratified & ExtraTrees A0 & 0.5095 & 0.6272 \\
state-stratified & ExtraTrees A3 & 0.5035 & 0.6360 \\
state-stratified & ExtraTrees A7 & \textbf{0.4966} & 0.6498 \\
state-stratified & Ridge A3 & 0.5298 & 0.6025 \\
state-stratified & Ridge A7 & 0.6515 & 0.3964 \\
state-stratified & DummyMean & 0.8506 & -0.0179 \\
\midrule
municipality-grouped & Ridge A3 & \textbf{0.7426} & -5.1998 \\
municipality-grouped & Ridge A6 & 0.8178 & -4.0401 \\
municipality-grouped & ExtraTrees A7 & 0.8916 & -6.4127 \\
municipality-grouped & DummyMean & 1.0066 & -6.1150 \\
\bottomrule
\end{tabular}
\end{table}

Los $R^2$ muy negativos de algunos folds grouped no implican necesariamente una implementación
incorrecta. Con folds pequeños y distribuciones de $y$ localmente estrechas, el denominador
de $R^2$ puede ser pequeño; un modelo con error absoluto moderado puede quedar muy por debajo
del predictor de la media del fold. Por eso RMSE es la métrica central.

\subsection{Out-of-fold como objeto de primera clase}

No basta con guardar una media de scores. Para cada modelo se conservan predicciones OOF:
\[
\widehat y_i^{\mathrm{OOF}}
=
f_{\widehat\theta_{-f(i)}}(x_i),
\]
donde $f(i)$ es el fold que contiene a $i$. Esto permite posteriormente:
\begin{itemize}
\item computar métricas pooled sobre las 138 parcelas;
\item analizar residuos espacialmente;
\item formar ensambles sin usar predicción in-sample;
\item detectar si un fold o municipio domina una métrica.
\end{itemize}
La persistencia de OOF fue crucial en 04A, donde los residuos de ensambles convencionales se
analizaron con Moran's $I$.

\subsection{Hallazgo metodológico}

El checkpoint revela una tensión importante: bajo state-stratified, ExtraTrees con competition
completo se acerca a 0.50 RMSE; bajo municipality-grouped, los mismos métodos pueden superar
0.85--0.90. Esta brecha es mucho mayor que las diferencias de milésimas entre
hiperparámetros.

La consecuencia fue priorizar:
\begin{enumerate}
\item representaciones compactas con estructura agronómica;
\item evaluación grouped como señal de robustez;
\item selección fold-local;
\item más adelante, métodos explícitamente espaciales/locales.
\end{enumerate}

Checkpoint~02 no seleccionó el modelo final. Congeló el terreno experimental sobre el que
Checkpoint~03 pudo comparar ingeniería de representaciones sin cambiar simultáneamente folds,
preprocessing y familia de estimador.

\chapter{Checkpoint 03: Ingeniería de representaciones y primera entrega de modelado}
\label{chap:checkpoint03}
\section{03A. Capa agronómica no lineal}
\label{sec:checkpoint03a}
\subsection{Motivación: no linealidad con semántica agronómica}

El Checkpoint~03A introduce una capa aditiva de 350 variables derivadas sin utilizar
\artifact{RENDIMIENTO_T_HA}. El objetivo no es reemplazar Feature Table v1, sino codificar
hipótesis no lineales interpretables antes de abrir una búsqueda de modelos. Con aproximadamente
1,400 predictores base, una expansión polinómica completa de segundo orden generaría del orden
de $10^6$ interacciones; esa estrategia es inadecuada para $n=138$. En su lugar, se construyen
familias pequeñas cuya razón agronómica puede enunciarse antes de observar $y$.

El build final contiene 197 filas y 350 features: 123 clean y 227 competition. No hay missing
ni infinitos. Tres variables de calor resultaron constantes y se eliminaron automáticamente.
Las familias retenidas son:
\[
\begin{array}{lr}
\text{fenología} & 162\\
\text{anomalías fenológicas} & 41\\
\text{agua/productividad} & 40\\
\text{térmicas} & 36\\
\text{cross-domain} & 19\\
\text{perfil de suelo} & 18\\
\text{timing de agua} & 10\\
\text{sensor agreement} & 9\\
\text{interacciones de suelo} & 8\\
\text{bases no lineales} & 7
\end{array}
\]

\subsection{Fenología y forma estacional}

Para una trayectoria mensual $x_t$, $t\in\{4,\ldots,10\}$, se calculan media estacional,
amplitud, AUC, mes de máximo, medias temprana/tardía, cambio junio--agosto y pendientes de
green-up/senescence. Por ejemplo,
\begin{align}
\bar x &= \frac{1}{T}\sum_{t=1}^{T}x_t,\\
A(x) &= \max_t x_t-\min_t x_t,\\
AUC(x) &\approx
\sum_{t=1}^{T-1}\frac{x_t+x_{t+1}}{2}
(m_{t+1}-m_t),\\
t_{\max} &= \arg\max_t x_t.
\end{align}

Las pendientes son coeficientes de regresiones lineales locales sobre bloques de meses:
\[
x_t=\alpha+\beta t+\varepsilon_t.
\]
No se afirma que la fenología real sea lineal; $\beta$ funciona como descriptor compacto de
dirección de cambio. La literatura de cebada y estimación de rendimiento por teledetección
motiva conservar timing y forma en lugar de depender sólo de una media
\citep{Chanev2025,Sharifi2020,Mirosavljevic2018}.

\subsection{Anomalías 2025 frente a historia}

Una parcela puede tener productividad estructuralmente alta o baja y, al mismo tiempo, mostrar
una condición 2025 inusual respecto a su propia historia. La capa separa ambas dimensiones
mediante diferencias y razones:
\begin{align}
\Delta\bar x_i
&=
\bar x_{i,2025}-\bar x_{i,\mathrm{hist}},\\
R_{x,i}
&=
\frac{\bar x_{i,2025}}
{\bar x_{i,\mathrm{hist}}+\epsilon},\\
\Delta AUC_i
&=
AUC_{i,2025}-AUC_{i,\mathrm{hist}}.
\end{align}

También se incluyen delta de agosto, delta de late-season y desplazamiento del mes de pico.
Para NDWI se evita deliberadamente una razón de medias, porque el denominador puede cruzar
cero al tratarse de un índice signed. Esta es una decisión numérica y semántica: una feature
de gran magnitud causada por un denominador cercano a cero no debe confundirse con una señal
agronómica extrema.

\subsection{Concordancia inter-sensor}

Sentinel-2 y Planet 2025 se comparan para NDVI, EVI y LAI mediante diferencia de media,
razón de media y RMSE mensual:
\begin{equation}
\RMSE_{\mathrm{sensor},i}
=
\sqrt{
\frac{1}{T}\sum_t
\left(x^{S2}_{i,t}-x^{Planet}_{i,t}\right)^2
}.
\label{eq:sensor-rmse}
\end{equation}
La feature no presupone equivalencia de nivel entre productos. Su interpretación es
diagnóstica: si dos sensores describen de manera muy distinta la misma parcela, esa
discordancia puede contener información de cobertura, procesamiento, nubosidad o dinámica
vegetal.

\subsection{Desarrollo térmico}

A partir de temperatura máxima y mínima se construye
\[
T_{\mathrm{mean},m}=\frac{T_{\max,m}+T_{\min,m}}{2},
\qquad
DTR_m=T_{\max,m}-T_{\min,m}.
\]
Se agregan proxies mensuales de growing degree days:
\begin{equation}
GDD_{\mathrm{proxy}}(T_b)
=
\sum_m d_m\max(T_{\mathrm{mean},m}-T_b,0),
\label{eq:gdd-proxy}
\end{equation}
para $T_b=0^\circ$C y $5^\circ$C. Se usa explícitamente la palabra proxy porque los inputs
climáticos son mensuales, no diarios. La acumulación térmica es una construcción habitual en
fenología de cereales \citep{McMaster1997,McMaster2005,Hajkova2019}, pero el valor exacto aquí
no debe confundirse con un GDD agronómico diario de alta resolución.

Se propuso también
\[
H_{25}
=
\sum_m d_m\max(T_{\mathrm{mean},m}-25,0).
\]
En la build canónica resultó constante tanto histórica como contemporáneamente y fue eliminado,
junto con la interacción que lo utilizaba. Este comportamiento es útil: el pipeline no fuerza
a conservar una hipótesis sólo porque conceptualmente sea razonable.

\subsection{Timing de precipitación}

La precipitación se separa en ventanas:
\begin{align}
P_{\mathrm{early}} &= P_{\mathrm{Apr}}+P_{\mathrm{May}}+P_{\mathrm{Jun}},\\
P_{\mathrm{mid}} &= P_{\mathrm{Jun}}+P_{\mathrm{Jul}}+P_{\mathrm{Aug}},\\
P_{\mathrm{late}} &= P_{\mathrm{Aug}}+P_{\mathrm{Sep}}+P_{\mathrm{Oct}}.
\end{align}
También se calculan coeficiente de variación y fracción tardía. La motivación es que la
respuesta del cultivo a déficit hídrico depende de la etapa de desarrollo
\citep{Bello2022,Mogensen1980,Zhang2015}. El total estacional no contiene por sí solo
información sobre cuándo ocurrió el agua.

\subsection{Productividad hídrica y WaPOR}

AETI y NPP se combinan en:
\begin{align}
WP_{\mathrm{NPP}}
&=
\frac{NPP}{AETI+\epsilon},\\
R_{ET/P}
&=
\frac{AETI}{P+\epsilon},\\
WB
&=
P-AETI.
\end{align}
La primera es un proxy de productividad de biomasa por agua observada desde teledetección, no
una definición de water-use efficiency de grano. De forma similar, $AETI/P$ no es un índice de
sequía: el cociente puede estar influido por almacenamiento de agua en suelo, riego u otras
fuentes. La nomenclatura deliberadamente evita atribuir causalidad fisiológica a una
transformación predictiva \citep{Dhonthi2024}.

Se añaden interacciones como
\[
NPP\times NDVI,\qquad LAI\times AETI,
\]
para permitir a modelos lineales expresar que la biomasa y el uso de agua pueden tener
efectos dependientes del estado del dosel.

\subsection{Perfil edáfico y bases no lineales}

Para cada propiedad $p$ de SoilGrids se construyen contraste vertical y razón:
\begin{align}
G_p &= p_{0-30}-p_{30-60},\\
R_p &= \frac{p_{0-30}}{p_{30-60}+\epsilon}.
\end{align}
Adicionalmente se crean interacciones compactas: SOC$\times$CEC,
SOC$\times$N, clay$\times$SOC, clay$\times$CEC, bulk-density$\times$clay,
clay/sand, clay+silt y $pH^2$. La forma cuadrática de pH no fija un óptimo; simplemente
permite que un modelo lineal represente curvatura junto con el término lineal. La respuesta de
cebada a química y cationes de suelo puede ser no lineal \citep{Holland2021}.

Variables positivas y sesgadas reciben una base
\[
z=\log(1+x),
\]
principalmente para estabilizar escalas en modelos lineales.

\subsection{Interacciones cross-domain}

Se crean hipótesis entre dominios, entre ellas:
\[
Y_{\mathrm{SIAP,hist}}\times\Delta NDVI_{2025},
\qquad
Y_{\mathrm{SIAP,hist}}\times NPP_{2025},
\]
así como slope$\times$precipitación, elevación$\times$temperatura,
clay$\times$precipitación, SOC$\times$precipitación, N de suelo$\times$NDVI,
NDVI$\times$NDWI y EVI$\times$LAI. El objetivo es proporcionar bases físicamente
interpretables sin enumerar todas las interacciones posibles.

\subsection{Clases de evidencia}

El manifest asigna a cada feature una clase:
\begin{description}
\item[literature-backed] la relación general está motivada por literatura;
\item[mechanistic proxy] la transformación tiene interpretación física/agronómica clara,
pero la fórmula exacta es un proxy del proyecto;
\item[experimental] base matemática target-free que debe demostrar valor predictivo.
\end{description}

Esta clasificación impide escribir retrospectivamente que una interacción que funcionó en CV
es una ley agronómica. La validez predictiva se decide en Checkpoint~03C; la validez
mecanicista no se infiere del score.

\section{03B. Capa empírica y descubrimiento fold-local de expresiones}
\label{sec:checkpoint03b}
\subsection{Objetivo: estructura empírica sin target}

Checkpoint~03B responde una pregunta distinta de 03A. En lugar de codificar primero una
hipótesis agronómica, construye descriptores matemáticos de la geometría observada de $X$:
curvatura temporal, persistencia, cambio relativo robusto, concordancia de forma entre
sensores y posición de 2025 respecto al historial corto.

El build canónico retiene 335 features en 197 parcelas, sin missing ni infinitos:
91 clean y 244 competition. Las familias son:
\[
162\ \text{temporal shape}
+42\ \text{symmetric change}
+41\ \text{short-baseline condition}
+72\ \text{aggregate geometry}
+18\ \text{sensor shape similarity}.
\]
Sólo se elimina una columna constante. El target no participa en la construcción de esta
tabla.

\subsection{Geometría temporal}

Dada una serie $x_1,\ldots,x_T$, se definen diferencias locales
\[
s_t=\frac{x_t-x_{t-1}}{m_t-m_{t-1}}.
\]
A partir de ellas se construyen:

\paragraph{Variación total}
\begin{equation}
TV(x)=\sum_{t=2}^{T}|x_t-x_{t-1}|.
\label{eq:tv}
\end{equation}

\paragraph{Rugosidad}
\begin{equation}
R(x)=\frac{1}{T-2}\sum_t|s_t-s_{t-1}|.
\label{eq:roughness}
\end{equation}

\paragraph{Energía de curvatura}
\begin{equation}
C(x)=\frac{1}{T-2}\sum_t(s_t-s_{t-1})^2.
\label{eq:curvature}
\end{equation}

La variación total mide movimiento acumulado, mientras que rugosidad y energía de curvatura
penalizan cambios en la pendiente. Dos parcelas pueden tener la misma amplitud y AUC, pero
trayectorias muy distintas en estas métricas.

También se calculan autocorrelación lag-1, $R^2$ de tendencia lineal, fracción de cambios de
signo en las pendientes, sharpness del máximo, centro de masa temporal y entropía de forma.
Para el centro de masa se normaliza primero:
\[
w_t=\frac{x_t-\min x}{\max x-\min x},
\qquad
M(x)=\frac{\sum_t m_t w_t}{\sum_t w_t}.
\]
Para la entropía, si $p_t=w_t/\sum_s w_s$,
\[
H(x)=-\frac{\sum_t p_t\log p_t}{\log T}.
\]
$H$ cercano a uno describe una señal temporal difusa; valores menores concentran masa en
pocos meses.

\subsection{Condición 2025 respecto a un historial corto}

La posición del valor 2025 dentro del rango histórico parcelario se define como
\begin{equation}
Q_{i,m}
=
\frac{x_{i,2025,m}-x_{i,\mathrm{hist,min}}}
{x_{i,\mathrm{hist,max}}-x_{i,\mathrm{hist,min}}+\epsilon}.
\label{eq:hist-position}
\end{equation}

Para NDVI se conserva una versión escalada y nombrada explícitamente
\artifact{vci_like_short}. Está inspirada en la lógica de índices de condición vegetal
\citep{Bokusheva2016,Serban2025}, pero no se presenta como VCI estándar: el historial sólo
abarca 2022--2024 y los extremos son resúmenes parcelarios, no una climatología larga.
Además, los valores no se recortan a $[0,1]$, de modo que una condición 2025 fuera del rango
histórico sigue siendo observable.

\subsection{Cambio simétrico estable}

Los índices espectrales signed pueden acercarse a cero; una razón ordinaria puede explotar.
Se introduce el contraste
\begin{equation}
D(a,b)=\frac{2(a-b)}{|a|+|b|+\epsilon}.
\label{eq:symmetric-change}
\end{equation}
Esta transformación es adimensional, simétrica respecto a la magnitud de ambos términos y
evita el comportamiento inestable de $a/b$ cuando $b\approx0$.

Para una secuencia de cambios mensuales se resumen media estacional, RMS, máximo absoluto,
media tardía, fracción de cambios positivos y fracción de inversiones de signo. Estos
descriptores son empíricos; no se etiquetan como índices agronómicos canónicos.

\subsection{Similitud temporal entre sensores}

Además del error de nivel de 03A, 03B compara \emph{forma} temporal Sentinel-2 versus Planet
mediante Pearson, Spearman, cosine similarity y RMSE normalizado por rango combinado. Sea
$\bm x$ una trayectoria S2 y $\bm z$ una Planet:
\[
\cos(\bm x,\bm z)
=
\frac{\bm x^\top\bm z}
{\|\bm x\|_2\|\bm z\|_2}.
\]
Las métricas no requieren que ambos sensores tengan calibración idéntica para detectar si
suben y bajan de manera concordante.

\subsection{Descubrimiento supervisado separado del dataset materializado}

La contribución metodológica más delicada de 03B es separar la capa $X$-only de un
\emph{miner} supervisado que nunca se materializa globalmente. Se configuran 24 primitivas.
Dentro de cada training fold:
\begin{enumerate}
\item se rankean primitivas usando únicamente $y$ visible en ese training fold;
\item se conservan las 12 primeras;
\item se generan expresiones de una familia predefinida: sumas, diferencias, productos,
safe ratios, reverse ratios y symmetric change;
\item se rankean por asociación Spearman dentro del training fold;
\item se conservan hasta 20;
\item esas expresiones se transforman sobre train/validation sin volver a usar $y$ de
validación.
\end{enumerate}

Formalmente, si $A$ denota el algoritmo de descubrimiento y $f$ un fold,
\[
\mathcal E_f=A(X_{\mathrm{train}(f)},y_{\mathrm{train}(f)}),
\]
y \emph{no}
\[
\mathcal E=A(X_{\Train},y_{\Train})
\]
reutilizado luego en los folds. La segunda construcción filtraría información de los holdouts
hacia la representación.

\subsection{Auditoría de recurrencia}

La auditoría de estabilidad usa 138 parcelas, 24 primitivas, 12 primitivas por fit y 20
expresiones por fit. A través de los diez folds de los dos protocolos se seleccionaron 88
expresiones únicas. El objetivo no era estimar performance, sino identificar si el miner
producía estructuras recurrentes.

Siete expresiones aparecieron al menos en tres de cinco folds bajo \emph{ambos} protocolos.
Sus componentes se concentran en:
\begin{itemize}
\item tendencia y media reciente de rendimiento municipal histórico SIAP;
\item precipitación estacional 2025;
\item GDD proxy base $5^\circ$C;
\item carbono orgánico de suelo 0--30 cm.
\end{itemize}

Ejemplos estructurales:
\[
\frac{P_{2025}}
{|\mathrm{SIAPtrend}|+\epsilon},
\qquad
\frac{SOC_{0-30}}
{|P_{2025}|+\epsilon},
\]
y
\[
\frac{2(\mathrm{SIAPtrend}-P_{2025})}
{|\mathrm{SIAPtrend}|+|P_{2025}|+\epsilon}.
\]

La alta correlación Spearman de entrenamiento de estas fórmulas no es evidencia de
generalización. Sólo 03C, refitando el miner dentro de cada outer fold, puede responder si la
familia de descubrimiento agrega valor predictivo.

\subsection{Resultado conceptual}

03B evita dos errores frecuentes en feature engineering small-$n$:
\begin{enumerate}
\item generar miles de combinaciones globales y escoger las que correlacionan con todo
$y_{\Train}$ antes de CV;
\item confundir estabilidad de una asociación dentro de train con performance out-of-sample.
\end{enumerate}
La capa determinista queda disponible para modelos y embeddings; el descubrimiento supervisado
queda encapsulado como transformer fold-local.

\section{03C.1. Benchmark controlado de representaciones}
\label{sec:checkpoint03c1}
\subsection{Pregunta experimental}

Con Feature Table v1, la capa agronómica y la capa empírica disponibles, Checkpoint~03C.1
separa el efecto de la \emph{representación} del efecto de la familia de modelo. Para hacerlo,
usa exactamente los mismos dos estimadores de referencia ---Ridge10 y ExtraTrees--- sobre
siete representaciones por track y conserva los folds de Checkpoint~02.

El diseño evita una comparación confusa del tipo representación A + modelo muy ajustado contra
representación B + modelo distinto. Aquí la pregunta es:
\[
\Delta_{\mathrm{representation}}
=
\RMSE(M,X^{(r_2)})-\RMSE(M,X^{(r_1)})
\]
para un mismo modelo $M$ y un mismo protocolo.

\subsection{Representaciones evaluadas}

\begin{table}[H]
\centering
\caption{Representaciones de 03C.1.}
\label{tab:cp03c1-reps}
\small
\begin{tabular}{llrr}
\toprule
Track & Representación & $p$ & Discovery \\
\midrule
clean & B0 base & 692 & no \\
clean & B1 agronomic & 123 & no \\
clean & B2 empirical & 91 & no \\
clean & B3 base+agro & 815 & no \\
clean & B4 base+empirical & 783 & no \\
clean & B5 all & 906 & no \\
clean & B7 all+discovery & 906 & sí, hasta 20 \\
\midrule
competition & B0 base & 1,400 & no \\
competition & B1 agronomic & 350 & no \\
competition & B2 empirical & 335 & no \\
competition & B3 base+agro & 1,750 & no \\
competition & B4 base+empirical & 1,735 & no \\
competition & B5 all & 2,085 & no \\
competition & B7 all+discovery & 2,085 & sí, hasta 20 \\
\bottomrule
\end{tabular}
\end{table}

En B7 las columnas base son fijas, pero las expresiones supervisadas no lo son: cada outer
training fold descubre sus propias expresiones. El número 20 indica un máximo por fit, no un
set global de 20 fórmulas preseleccionadas.

\subsection{Métrica OOF versus promedio de folds}

03C.1 conserva dos resúmenes que no son idénticos cuando los folds tienen tamaños distintos.
Si $F$ es el conjunto de folds:
\[
\overline{\RMSE}_{\mathrm{fold}}
=
\frac{1}{|F|}
\sum_{f\in F}\RMSE_f,
\]
mientras
\[
\RMSE_{\mathrm{OOF}}
=
\sqrt{
\frac{1}{138}
\sum_{i=1}^{138}
(y_i-\hat y_i^{\mathrm{OOF}})^2
}.
\]
El primero pondera cada fold por igual; el segundo pondera cada parcela por igual. En
municipality-grouped, donde los folds pueden ser muy desiguales, reportar ambos evita esconder
la influencia de un municipio grande.

\subsection{Resultados en state-stratified}

Para competition + ExtraTrees, B7 obtiene OOF RMSE 0.4992, B0 0.5003, B1 0.5041 y B5
0.5045. Es decir, bajo este protocolo las diferencias entre representaciones ricas son
pequeñas; la capa empírica sola B2 es claramente peor, con 0.6470.

Con Ridge10 el comportamiento cambia: B0 alcanza 0.6534, B4 0.6587, B3 0.6644, B7 0.6679 y
B1 0.6809. La compactación agronómica que puede ser razonable para árboles no mejora
automáticamente un Ridge de alpha fijo en este track.

En clean, ExtraTrees B7 alcanza 0.4948 OOF y B3 0.5021, mientras Ridge sobre la capa
agronómica compacta B1 alcanza 0.5888 frente a 0.6718 de B0. Esto muestra que el valor de una
representación depende fuertemente de la inductive bias del modelo.

\subsection{Resultados en municipality-grouped}

Bajo el protocolo espacialmente más exigente, competition + ExtraTrees favorece B1 agronomic:
\[
\RMSE_{\mathrm{OOF}}=0.8590,
\]
seguido por B7=0.8675 y B0=0.8974. Para Ridge, B0 es claramente mejor:
0.8059, frente a B4=0.8464, B3=0.8676, B5=0.8976 y B7=0.9142.

En clean, ExtraTrees B7 llega a 0.8488 y B1 a 0.8762; Ridge B0 queda en 0.8741. Ninguna
representación produce una mejora universal. La diferencia entre protocolos sigue siendo más
grande que muchas diferencias internas de representación.

\subsection{Lección de 03C.1}

Tres conclusiones pasan a 03C.2:

\begin{enumerate}
\item \textbf{C0/base debe sobrevivir}: pese a la alta dimensionalidad, sigue siendo una
representación competitiva para modelos lineales latentes como PLS/Ridge.
\item \textbf{C1/agronomic merece un benchmark propio}: su compactación funciona bien con
árboles bajo el stress grouped y ofrece una representación de sólo 350 columnas.
\item \textbf{El bulk empírico no se añade por default}: B2 solo es débil y B5 puede
degradar algunos modelos; las primitivas empíricas más útiles se retienen de forma controlada
para discovery.
\end{enumerate}

Por ello, 03C.2 reduce el espacio de representaciones en vez de continuar acumulando capas.
Este es un patrón deliberado del proyecto: cada checkpoint debe \emph{cerrar} grados de
libertad antes de abrir una búsqueda más costosa.

\section{03C.2. Benchmark anidado de familias de modelos}
\label{sec:checkpoint03c2}
\subsection{Diseño del benchmark anidado}

Checkpoint~03C.2 abandona el benchmark de estimadores fijos y permite tuning disciplinado.
El track activo pasa a ser exclusivamente competition. Se conservan cuatro representaciones:
\begin{align*}
C0 &: \text{base}, && p=1400,\\
C1 &: \text{agronomic}, && p=350,\\
C2 &: \text{base+agro+empirical+discovery}, && p=2085+20,\\
C3 &: \text{base+agro+support empírico+discovery}, && p=1754+20.
\end{align*}

C3 es particularmente importante: no arrastra las 335 features empíricas, pero conserva las
primitivas configuradas necesarias para que el miner pueda explorar relaciones tipo
condición relativa/shape. Así se distingue entre usar una capa empírica completa y usar sólo
su soporte para descubrimiento.

\subsection{Nested cross-validation}

Sea $f$ un outer fold. El conjunto outer-train se divide nuevamente en tres inner folds.
Para cada configuración $\lambda\in\Lambda$:
\[
\widehat R_f(\lambda)
=
\frac{1}{3}\sum_{g=1}^{3}
\RMSE\left(
y_{V_{fg}},
\hat f_{\lambda,T_{fg}}(X_{V_{fg}})
\right),
\]
donde cualquier imputación, escalado, PCA, PLS y descubrimiento supervisado se ajusta sólo en
$T_{fg}$. Se selecciona
\[
\hat\lambda_f
=
\arg\min_{\lambda\in\Lambda}\widehat R_f(\lambda),
\]
se refita el pipeline completo sobre el outer-train y sólo entonces se predice el outer
holdout.

El inner splitter refleja el propósito del outer protocol:
\begin{itemize}
\item outer state-stratified $\rightarrow$ inner state-stratified;
\item outer municipality-grouped $\rightarrow$ inner group-by-municipality.
\end{itemize}
Esto evita seleccionar hiperparámetros bajo una geometría de validación que no corresponde a
la evaluación exterior.

\subsection{Familias de modelos}

Se evaluaron Ridge, ElasticNet, ExtraTrees, CatBoost, PLS, PCA+RBF Kernel Ridge y
PCA+polynomial-degree-2 Kernel Ridge.

\paragraph{PLS.}
PLS busca componentes latentes maximizando covarianza con $y$, por lo que puede ser útil en
high-$p$/small-$n$ con covariables correlacionadas \citep{Wold2001}. Como es supervisado,
sus componentes se ajustan exclusivamente dentro de train.

\paragraph{CatBoost.}
El boosting de árboles se ejecutó en GPU y se restringió a grids pequeños
\citep{CatBoost2018}. El objetivo era comprobar si una no linealidad potente podía extraer
señal de la capa agronómica compacta sin un barrido masivo.

\paragraph{Kernels.}
Los kernels no se materializan como datasets porque dependen de preprocessing e
hiperparámetros. Para un RBF:
\[
K_\gamma(x,z)=\exp\{-\gamma\|x-z\|_2^2\},
\]
por lo que escalar $X$ fuera del fold alteraría directamente la geometría del modelo. El mismo
principio justifica mantener PCA dentro del pipeline.

\subsection{Ranking cross-protocol}

La tabla de robustez ordena por el peor OOF RMSE entre state y grouped, y después por la media.
Los tres finalistas individuales fueron:

\begin{table}[H]
\centering
\caption{Finalistas individuales de 03C.2.}
\label{tab:cp03c2-finalists}
\begin{tabular}{llrr}
\toprule
Representación & Modelo & state OOF & grouped OOF \\
\midrule
C0 base & PLS & 0.5339 & \textbf{0.7621} \\
C3 base+agro+discovery & Ridge & 0.5331 & 0.7650 \\
C1 agronomic & CatBoost & \textbf{0.5235} & 0.7709 \\
\bottomrule
\end{tabular}
\end{table}

El orden state por sí solo sería distinto. ExtraTrees C3 obtiene 0.4972, C2 0.4996 y C1
0.5024, mejores que los tres finalistas en state, pero sus grouped RMSE son 0.8698, 0.9052 y
0.8469 respectivamente. Por tanto, el proyecto no elige simplemente el mínimo state score:
se penaliza la fragilidad espacial.

\subsection{Protocol gap}

Definimos para diagnóstico
\[
G(M)=
\RMSE_{\mathrm{grouped}}(M)
-
\RMSE_{\mathrm{state}}(M).
\]
Para C0+PLS,
\[
G\approx0.2282;
\]
para C1+CatBoost,
\[
G\approx0.2474.
\]
Incluso los mejores finalistas muestran una brecha sustancial. Esto anticipa una conclusión
fundamental de 04A: el rendimiento y los residuos tienen estructura espacial que los modelos
globales no representan completamente bajo extrapolación por municipio.

\subsection{Fracaso de PCA+RBF y PCA+Poly2}

Los kernels no lineales presentan errores muy superiores: PCA+RBF alcanza aproximadamente
3.24--4.05 RMSE, y PCA+Poly2 entre 0.72 y más de 2 dependiendo de representación/protocolo.
Esto no significa que los kernels sean universalmente inadecuados. Significa que, con sólo 138
etiquetas, el grid y preprocessing definidos no produjeron una geometría robusta. El proyecto
no siguió gastando compute en una familia que ya mostraba una brecha enorme respecto a PLS,
Ridge y árboles.

\subsection{Por qué no se hizo una búsqueda gigantesca}

Con $n=138$, aumentar el número de configuraciones evaluadas sobre los mismos folds incrementa
la posibilidad de seleccionar ruido de validación. El compute disponible no es la restricción
principal; la restricción es la cantidad de información independiente. Por ello 03C.2 utiliza
candidate grids deterministas y relativamente pequeños.

Esta política se mantiene posteriormente en 04C.1: aun disponiendo de GPU, sólo se reevalúan
tres configuraciones CatBoost heredadas cuando ya existe una hipótesis concreta de diversidad.

\subsection{Primera lectura científica}

03C.2 demuestra que:
\begin{enumerate}
\item una representación simple C0 + PLS puede ser más robusta espacialmente que ensembles
de árboles con mejor state RMSE;
\item C1 agronomic contiene suficiente estructura para que CatBoost sea competitivo;
\item la capa completa C2 no domina consistentemente a representaciones más pequeñas;
\item el modelo global todavía sufre cuando se retienen municipios enteros.
\end{enumerate}

El checkpoint constituye el núcleo de la \emph{First Modeling Delivery}, pero todavía no usa
la geometría específica de las 59 parcelas como parte explícita del proceso de selección.

\section{03C.3. Ensambles finalistas de la primera entrega}
\label{sec:checkpoint03c3}
\subsection{Propósito del ensamble}

Después de 03C.2 se congelan tres finalistas individuales:
\[
F_1=\text{C0+PLS},\qquad
F_2=\text{C3+Ridge},\qquad
F_3=\text{C1+CatBoost}.
\]
Checkpoint~03C.3 no abre tuning adicional. Su objetivo es comprobar si combinaciones
deterministas y de pesos fijos reducen error OOF al promediar inductive biases distintos.

Los ensambles son:
\begin{align}
E_{12}&=\frac{1}{2}(F_1+F_2),\\
E_{13}&=\frac{1}{2}(F_1+F_3),\\
E_{23}&=\frac{1}{2}(F_2+F_3),\\
E_{123}&=\frac{1}{3}(F_1+F_2+F_3).
\end{align}

Los pesos no se optimizan contra OOF; de esta forma, 03C.3 mide únicamente el beneficio de
promediado y diversidad entre finalistas.

\subsection{Resultados}

\begin{table}[H]
\centering
\caption{Primera entrega de modelos: RMSE OOF por protocolo.}
\label{tab:cp03c3-results}
\begin{tabular}{lrrr}
\toprule
Método & state & municipality-grouped & peor protocolo \\
\midrule
E123 & 0.5121 & \textbf{0.7481} & 0.7481 \\
E13 & \textbf{0.5082} & 0.7488 & 0.7488 \\
E23 & 0.5146 & 0.7544 & 0.7544 \\
E12 & 0.5270 & 0.7546 & 0.7546 \\
F1 C0+PLS & 0.5339 & 0.7621 & 0.7621 \\
F2 C3+Ridge & 0.5331 & 0.7650 & 0.7650 \\
F3 C1+CatBoost & 0.5235 & 0.7709 & 0.7709 \\
\bottomrule
\end{tabular}
\end{table}

El ensamble reduce error respecto a cada componente. E13 es ligeramente mejor en state,
mientras E123 es ligeramente mejor en grouped. La diferencia grouped de $0.0007$ entre ambos
es demasiado pequeña para declarar un ganador sustantivo.

\subsection{Interpretación de diversidad}

Para dos predictores $\hat y^{(a)}$ y $\hat y^{(b)}$, el MSE del promedio es
\[
\MSE\left(\frac{\hat y^{(a)}+\hat y^{(b)}}{2}\right)
=
\frac{1}{4}
\left[
\MSE_a+\MSE_b
+2\,\mathbb E(e_a e_b)
\right],
\]
ignorando aquí términos de notación asociados a sesgo. El promedio sólo puede aportar una
reducción importante si los errores no están perfectamente correlacionados. PLS, Ridge y
CatBoost constituyen inductive biases diferentes, por lo que el resultado de 03C.3 es
coherente con diversidad parcial.

Sin embargo, el beneficio del ensamble no resuelve la fragilidad espacial: grouped permanece
cerca de 0.75 RMSE. Esto sugiere que la diversidad entre \emph{modelos globales} no sustituye
una representación explícita de localidad.

\subsection{Cierre de la First Modeling Delivery}

Al terminar 03C.3, el proyecto dispone de:
\begin{itemize}
\item data contract y Feature Table reproducibles;
\item capas agronómica y empírica;
\item discovery supervisado fold-local;
\item dos protocolos de CV congelados;
\item un benchmark anidado de siete familias de modelos;
\item una pequeña familia de ensambles sin tuning de pesos.
\end{itemize}

La mejor evidencia convencional es aproximadamente 0.508--0.512 state y 0.748 grouped. Este
resultado es suficientemente fuerte como baseline, pero revela la limitación central que
motiva Checkpoint~04: el conjunto real de 59 parcelas es conocido y su relación espacial con
las 138 etiquetas no está siendo utilizada explícitamente.

\begin{decision}
Checkpoint~03 se congela como evidencia histórica. Sus folds y OOF siguen siendo útiles para
diagnóstico, pero dejan de ser la única doctrina de selección. Checkpoint~04 define validación
pseudo-competition que reproduce la información realmente disponible al predecir las 59
parcelas.
\end{decision}

\chapter{Checkpoint 04: Modelado transductivo para los 59 objetivos fijos}
\label{chap:checkpoint04}
\section{04A. Topología del conjunto objetivo, similitud y autocorrelación espacial}
\label{sec:checkpoint04a}
\subsection{Reformulación transductiva}

Checkpoint~04 cambia la pregunta de modelado. Sea $\Train$ el conjunto etiquetado y
$\Target$ el conjunto fijo de 59 parcelas objetivo. La estimación final tiene forma
\[
\widehat{\bm y}_{\Target}
=
\mathcal A(X_{\Train},y_{\Train},X_{\Target}),
\]
donde $\mathcal A$ puede explotar la geometría conjunta de todos los $X$. Esto corresponde a
aprendizaje transductivo/semi-supervisado en el sentido de que el conjunto no etiquetado
específico participa en la representación \citep{Chapelle2006,Belkin2006}.

El primer paso no fue entrenar un modelo más complejo. Fue preguntar si la cercanía entre
parcelas realmente transfiere rendimiento y si los 59 objetivos forman un dominio distinto.

\subsection{Representaciones para topología}

Se construyeron tres espacios:
\begin{description}
\item[C0 base] Feature Table competition;
\item[C1 agronomic] las 350 features agronómicas;
\item[C4 all deterministic] base + agronomic + empirical, excluyendo expresiones
supervisadas descubiertas con $y$.
\end{description}

Para distancias multivariadas se imputa/estandariza usando exclusivamente $X$ y se obtiene un
embedding PCA de 20 componentes sobre las 197 parcelas. Esta operación es transductiva pero
$X$-only:
\[
Z=\mathrm{PCA}_{20}\bigl(\mathrm{scale}(X_{\Train}\cup X_{\Target})\bigr).
\]
No se usa $y$ para orientar los componentes.

\subsection{Clasificación adversarial train--target}

Se define una etiqueta auxiliar
\[
d_i=
\begin{cases}
0,&i\in\Train,\\
1,&i\in\Target.
\end{cases}
\]
Un clasificador cross-validated intenta distinguir ambos grupos a partir de C4/PCA. Si su AUC
estuviera muy por encima de 0.5, habría evidencia de covariate shift multivariado global. El
resultado fue
\[
\mathrm{AUC}=0.4892.
\]
Por tanto, los 59 objetivos no constituyen una población globalmente separable en la
representación examinada.

Este resultado no contradice que algunas parcelas individuales estén en bajo soporte. AUC
global cercano a azar implica solapamiento agregado, no que cada punto objetivo tenga un buen
vecino etiquetado.

\safefigure{../../reports/checkpoint_04a/figures/adversarial_roc.png}
{ROC de la clasificación adversarial train frente a target.}
{fig:cp04a-adversarial}

\subsection{¿Qué distancia transfiere rendimiento?}

Sobre las 138 etiquetas se construyeron todos los
\[
{138\choose2}=9{,}453
\]
pares. Para cada par $(i,j)$ se calcula
\[
\Delta y_{ij}=|y_i-y_j|
\]
y diversas métricas de disimilitud/similitud. La asociación se resume con Spearman, lo que
evita asumir linealidad entre distancia y diferencia de rendimiento.

\begin{table}[H]
\centering
\caption{Asociación entre similitud de covariables y diferencia absoluta de rendimiento.}
\label{tab:cp04a-transfer}
\begin{tabular}{lr}
\toprule
Métrica & Spearman con $|y_i-y_j|$ \\
\midrule
distancia geográfica & \textbf{0.4882} \\
distancia PCA C1 agronomic & 0.2200 \\
distancia PCA C4 deterministic & 0.1833 \\
distancia PCA C0 base & 0.1673 \\
DTW temporal & 0.0763 \\
Pearson temporal & -0.0503 \\
best-lag temporal correlation & -0.0482 \\
\bottomrule
\end{tabular}
\end{table}

Para distancias, una asociación positiva es la dirección útil: parcelas más lejanas tienden a
diferir más en rendimiento. Para similitudes tipo correlación, la dirección útil es negativa.
La geografía domina claramente a la similitud temporal mensual aislada.

\safefigure{../../reports/checkpoint_04a/figures/similarity_vs_yield_geo.png}
{Distancia geográfica frente a diferencia absoluta de rendimiento para pares etiquetados.}
{fig:cp04a-geo-transfer}

\subsection{Similitud temporal}

Se analizaron 11 trayectorias mensuales abril--octubre 2025: Sentinel-2 NDVI/EVI/LAI/FAPAR/
NDWI/MSI, Landsat VI6T y Planet NDVI/EVI/LAI/MSAVI. Para dos secuencias se calcularon
Pearson, Spearman, mejor correlación con lags de hasta dos meses y una distancia DTW
normalizada.

El resultado no es que la temporalidad carezca de información. Algunas parcelas tienen
análogos temporales convincentes. El hallazgo es más específico: \emph{como regla de
transferencia de rendimiento por sí sola}, la similitud mensual es mucho más débil que la
geografía. Por ello se evita construir el predictor final como nearest-neighbor temporal.

\subsection{Autocorrelación espacial de rendimiento}

La evidencia más fuerte de 04A proviene de Moran's $I$ \citep{Moran1950}. Para una matriz de
pesos espaciales $W=(w_{ij})$ y $z_i=y_i-\bar y$:
\begin{equation}
I
=
\frac{n}{\sum_{i,j}w_{ij}}
\frac{\sum_{i,j}w_{ij}z_i z_j}
{\sum_i z_i^2}.
\label{eq:moran}
\end{equation}

Sobre un grafo de cinco vecinos, el rendimiento observado produce
\[
I=0.6797,\qquad p_{\mathrm{perm}}=0.001,
\qquad \mathbb E_0[I]\approx-0.0073.
\]
Existe, por tanto, autocorrelación espacial positiva intensa.

\safefigure{../../reports/checkpoint_04a/figures/moran_yield.png}
{Moran scatter del rendimiento observado bajo el grafo espacial de cinco vecinos.}
{fig:cp04a-moran}

\subsection{Autocorrelación residual y dependencia del protocolo}

Los residuos OOF de los ensambles de 03C.3 muestran una diferencia reveladora:
\begin{table}[H]
\centering
\caption{Moran's $I$ de residuos OOF de la primera entrega.}
\label{tab:cp04a-residual-moran}
\begin{tabular}{lrr}
\toprule
Residuo & $I$ & $p_{\mathrm{perm}}$ \\
\midrule
E13 state-stratified & -0.0470 & 0.314 \\
E123 state-stratified & -0.0375 & 0.429 \\
E13 municipality-grouped & 0.5589 & 0.001 \\
E123 municipality-grouped & 0.5483 & 0.001 \\
\bottomrule
\end{tabular}
\end{table}

Cuando entrenamiento contiene parcelas cercanas, el modelo global absorbe parte de la
estructura espacial y sus residuos dejan de autocorrelacionarse fuertemente. Cuando se
retienen municipios, permanece una gran estructura espacial sin explicar. Esta observación
justifica modelos locales/grafo y explica la enorme brecha grouped de Checkpoint~03.

\subsection{Soporte de las 59 parcelas}

El diagnóstico 04A combina geografía, distancia C4, temporalidad, probabilidad adversarial y
consistencia de vecinos etiquetados en un score descriptivo. La taxonomía inicial fue:
\[
5\ \mathrm{A\ high}
+2\ \mathrm{B\ feature}
+5\ \mathrm{C\ geographic}
+13\ \mathrm{D\ extrapolation}
+34\ \mathrm{E\ mixed}.
\]

Entre los objetivos de mayor soporte aparecen AGC\_088, AGC\_190, AGC\_126, AGC\_128 y
AGC\_086. Entre los de menor soporte aparecen AGC\_173, AGC\_029, AGC\_120, AGC\_183 y
AGC\_003.

El score no es una probabilidad calibrada ni una predicción de rendimiento. Su función es
responder: ¿qué tan representado está un objetivo por la nube etiquetada?

\safefigure{../../reports/checkpoint_04a/figures/target_support_ranking.png}
{Ranking descriptivo de soporte de las 59 parcelas objetivo.}
{fig:cp04a-support}

\subsection{Decisión de transición}

Los resultados descartan dos rutas como prioridad: domain adaptation global agresiva y
nearest-neighbor temporal puro. En cambio, sostienen:
\begin{enumerate}
\item validación que reproduzca la geometría de las 59 parcelas;
\item local Ridge con vecindarios geográficos;
\item grafos sobre las 197 parcelas;
\item blends global/local sólo si reducen pseudo-test RMSE;
\item uso del soporte como hipótesis que debe validarse, no como regla manual.
\end{enumerate}

El siguiente checkpoint convierte estas hipótesis descriptivas en performance out-of-sample.

\section{04B. Validación pseudo-competition transductiva}
\label{sec:checkpoint04b}
\subsection{Por qué CV convencional ya no es suficiente}

Una validación convencional toma una partición aleatoria o agrupada de las 138 etiquetas, pero
no intenta que el pseudo-test se parezca específicamente a las 59 parcelas de competencia.
En un problema de objetivo fijo, la pregunta relevante es condicional a $X_{\Target}$.

Checkpoint~04B introduce una pseudo-competición. Para una máscara
$\Target^\star\subset\Train$:
\begin{enumerate}
\item se oculta $y_i$ para $i\in\Target^\star$;
\item se mantiene visible $X_{\Target^\star}$;
\item toda transformación $X$-only puede usar las 197 covariables;
\item cualquier fit supervisado usa sólo
$\Train^\star=\Train\setminus\Target^\star$;
\item se producen predicciones para $\Target^\star$;
\item sólo después se revela $y_{\Target^\star}$ para calcular métricas.
\end{enumerate}

De esta forma la frontera de información de cada pseudo-split reproduce la frontera real:
\[
(X_{\Train^\star},y_{\Train^\star},X_{\Target^\star})
\quad\longrightarrow\quad
\widehat y_{\Target^\star}.
\]

\subsection{Tamaño de la pseudo-competición}

La competencia tiene una fracción objetivo
\[
\pi_U=\frac{59}{197}\approx0.2995.
\]
Aplicada al universo de 138 etiquetas:
\[
138\pi_U\approx41.3.
\]
Por eso cada pseudo-competición target-matched y state-random oculta 41 parcelas.

Se congelaron:
\begin{itemize}
\item 16 target-matched repeats;
\item 8 state-random repeats;
\item 5 folds state-stratified legados;
\item 5 folds municipality-grouped legados.
\end{itemize}
En total se evalúan 34 splits con propósitos distintos.

\subsection{Construcción target-matched}

Las máscaras target-matched no usan $y$. Entre múltiples candidatos se busca aproximar la
geometría del conjunto real en:
\begin{itemize}
\item composición state/municipality;
\item nearest geographic distance;
\item nearest agronomic distance;
\item temporal similarity;
\item adversarial target probability.
\end{itemize}

La calidad se resume, entre otros, por total variation de la distribución municipal y distancia
entre perfiles estandarizados. El procedimiento logra:
\[
TV_{\mathrm{municipio}}:
0.0547\ \text{target-matched}
\quad\text{vs}\quad
0.0889\ \text{state-random},
\]
y
\[
d_{\mathrm{perfil}}:
0.5369\ \text{target-matched}
\quad\text{vs}\quad
0.9227\ \text{state-random}.
\]
Esto verifica que el protocolo principal se parece más a los 59 objetivos en $X$, sin
conocer sus $y$.

\safefigure{../../reports/checkpoint_04b/figures/split_match_quality.png}
{Calidad de matching de las pseudo-competiciones respecto al conjunto objetivo real.}
{fig:cp04b-match}

\subsection{Modelos candidatos}

04B mezcla baselines globales y métodos explícitamente locales.

\paragraph{Geographic kNN.}
Para una consulta $q$ y sus $k$ vecinos visibles $\mathcal N_k(q)$:
\[
\hat y_q
=
\frac{\sum_{i\in\mathcal N_k(q)}
(d_{qi}+\epsilon)^{-p}y_i}
{\sum_{i\in\mathcal N_k(q)}
(d_{qi}+\epsilon)^{-p}}.
\]

\paragraph{Local Ridge.}
Para cada consulta se construye un training set diferente de $k$ vecinos y se resuelve
\begin{equation}
\hat\beta_q
=
\arg\min_{\beta}
\sum_{i\in\mathcal N_k(q)}
w_{qi}
(y_i-x_i^\top\beta)^2
+\alpha\|\beta\|_2^2,
\qquad
w_{qi}=(d_{qi}+\epsilon)^{-1}.
\label{eq:local-ridge}
\end{equation}
La predicción es $\hat y_q=x_q^\top\hat\beta_q$. Este modelo es query-specific: dos objetivos
pueden usar subconjuntos de etiquetas diferentes.

\paragraph{Graph Laplacian.}
A partir de una matriz de distancias $D_X$ se crea un kNN simétrico con pesos RBF:
\[
w_{ij}
=
\exp\left[
-\frac12\left(\frac{d_{ij}}{\sigma}\right)^2
\right],
\]
donde $\sigma$ es la mediana de distancias de aristas positivas. Sea $A=(w_{ij})$,
$D=\diag(A\bm 1)$ y $L=D-A$. Con máscara $M$ de nodos etiquetados se resuelve
\begin{equation}
(M+\lambda L+\epsilon I)\bm f
=
M(\bm y-\bar y\bm 1).
\label{eq:graph-laplacian}
\end{equation}
La predicción en nodos query es $\bar y+f_q$. Esta formulación propaga una función suave sobre
la geometría $X$ \citep{Zhou2004,Belkin2006}.

\paragraph{Distancias mixtas.}
Se evaluaron combinaciones normalizadas:
\[
d_{\mathrm{mix}}
=
a\,d_{\mathrm{geo}}
+b\,d_{\mathrm{agro}}
+c\,d_{\mathrm{temporal}},
\qquad a+b+c=1.
\]
El mejor resultado inicial no requiere una componente temporal dominante, consistente con 04A.

\subsection{Resultados target-matched}

\begin{table}[H]
\centering
\caption{Ranking principal de Checkpoint 04B.}
\label{tab:cp04b-primary}
\begin{tabular}{lrrrr}
\toprule
Método & RMSE medio & mediana & peor split & comentario \\
\midrule
LocalRidge $k=20,\alpha=100$ & \textbf{0.4946} & 0.4778 & 0.6927 & mejor media \\
Graph $k=8,\lambda=2$ & 0.4974 & 0.4878 & 0.6995 & robusto \\
LocalRidge $k=30,\alpha=100$ & 0.4976 & 0.4847 & 0.6950 & cercano \\
GeoKNN $k=10$ & 0.5021 & 0.4794 & 0.7137 & baseline local \\
Blend PLS8+Graph8 & 0.5093 & 0.5012 & 0.6968 & blend \\
GlobalPLS C0 4 & 0.5379 & -- & -- & global \\
\bottomrule
\end{tabular}
\end{table}

El mejor local mejora aproximadamente 8\% respecto a GlobalPLS C0 en el mismo protocolo
target-matched:
\[
\frac{0.5379-0.4946}{0.5379}\approx0.0805.
\]
Esto es una ganancia material comparada con las diferencias de milésimas entre variantes
globales de 03C.

\safefigure{../../reports/checkpoint_04b/figures/primary_method_ranking.png}
{Ranking de métodos bajo las 16 pseudo-competiciones target-matched.}
{fig:cp04b-ranking}

\subsection{Robustez del grafo}

El Graph $k=8,\lambda=2$ obtiene:
\[
\begin{array}{lr}
\text{target-matched} & 0.4974\\
\text{state-random} & 0.4655\\
\text{state-stratified} & 0.4779\\
\text{municipality-grouped} & 0.5466.
\end{array}
\]
Su peor familia, 0.5466, es muy inferior al grouped 0.7481 de E123. No es una comparación
perfectamente idéntica porque 04B permite topología transductiva $X$-only; precisamente ése es
el punto: la estructura del problema permite un estimador más apropiado.

LocalRidge lidera en target-matched pero se deteriora mucho más al retirar municipios enteros.
La tensión entre Local y Graph se conserva hasta 04F: Local optimiza la simulación más cercana
a los 59 objetivos, Graph funciona como hedge de robustez espacial.

\subsection{Routing por soporte: primera hipótesis y limitación}

04B agrupó pseudo-targets en tiers low/mid/high según un score estrictamente $X$-only
recalculado dentro del split. En la tabla agregada:
\[
\begin{array}{lll}
\text{low}: & \text{LocalRidge30}, & \RMSE=0.4174,\\
\text{mid}: & \text{GeoKNN10}, & \RMSE=0.5645,\\
\text{high}: & \text{LocalRidge30}, & \RMSE=0.3198.
\end{array}
\]
Esto generó una propuesta inicial para las 59 parcelas: 32 GeoKNN10 y 27 LocalRidge30.

Sin embargo, la misma evidencia pseudo-test se utilizaba para identificar el ganador del tier
y medirlo. Por tanto, el routing es exploratorio. La solución correcta es seleccionar reglas
en splits distintos de aquellos donde se evalúan; 04D.1 introduce precisamente un
leave-one-split-out.

\subsection{Resultado de diseño}

04B cambia el centro del proyecto. Ya no es razonable invertir la mayor parte del esfuerzo en
modelos globales. La evidencia indica:
\begin{enumerate}
\item la geografía transfiere rendimiento;
\item LocalRidge puede superar el mejor anchor global en la geometría target-matched;
\item Graph ofrece una robustez extraordinariamente mejor al extrapolar espacialmente;
\item cualquier routing o tuning adicional debe reutilizar exactamente estas máscaras
congeladas para evitar mover simultáneamente objetivo y método.
\end{enumerate}

El siguiente gran bloque, 04D.1, refina sólo las familias locales/grafo que ya sobrevivieron
esta prueba.

\section{04C.1. Anchor global CatBoost enfocado}
\label{sec:checkpoint04c1}
\subsection{Motivación y posición cronológica}

Checkpoint~04C.1 se ejecutó después de haber observado la superioridad de métodos
locales/grafo y después de cerrar la hipótesis SIAP. Se mantiene la etiqueta 04C porque
representa la rama de modelos globales de mayor compute prevista desde el diseño de
Checkpoint~04. La pregunta no es volver a buscar el mejor modelo global, sino medir si
CatBoost aporta un error suficientemente diferente para mejorar un predictor local mediante
un peso pequeño.

La representación se restringe a C1 agronomic, con 350 features, porque 03C.2 ya había
identificado C1+CatBoost como el anchor boosting global más relevante. Reabrir cientos de
configuraciones sobre los mismos 138 labels habría aumentado el riesgo de overfitting de
selección sin aportar información independiente.

\subsection{Candidate set congelado}

Se reutilizan tres configuraciones presentes en el grid histórico de 03C.2:
\begin{table}[H]
\centering
\caption{Candidatos CatBoost de 04C.1.}
\label{tab:cp04c-cat}
\begin{tabular}{rrrr}
\toprule
depth & learning rate & $L_2$ leaf reg & iterations \\
\midrule
4 & 0.03 & 3 & 400 \\
4 & 0.06 & 6 & 400 \\
5 & 0.06 & 3 & 400 \\
\bottomrule
\end{tabular}
\end{table}

Cada configuración se entrena con seeds 42, 314 y 2718; la predicción usada para evaluación
es
\[
\hat y_q^{CB}
=
\frac{1}{3}
\sum_{s\in\{42,314,2718\}}
\hat y_{q,s}^{CB}.
\]
La ejecución se hace en GPU; no existe fallback automático a CPU. Esta decisión es de
reproducibilidad operacional, no estadística.

\subsection{Blends predefinidos}

Para el candidato estable depth 4, learning rate 0.03, $L_2=3$, se definen pesos pequeños
\[
\hat y^{(w)}
=
(1-w)\hat y^{Local}
+w\hat y^{CB},
\qquad
w\in\{0.10,0.20,0.30\}.
\]
También se construyen blends equivalentes con Graph04D como baseline. Los pesos son
deliberadamente discretos: dado que CatBoost standalone ya es más débil, la hipótesis plausible
es corrección/diversidad, no dominancia.

\subsection{Resultados target-matched}

\begin{table}[H]
\centering
\caption{Resultados principales de 04C.1.}
\label{tab:cp04c-results}
\small
\begin{tabular}{lrrr}
\toprule
Método & RMSE medio & pooled RMSE & pooled MAE \\
\midrule
Local + CatBoost 10\% & 0.487842 & 0.495915 & 0.342266 \\
Local04D & \textbf{0.487845} & \textbf{0.495741} & 0.342501 \\
Local + CatBoost 20\% & 0.488501 & 0.496728 & 0.342952 \\
Local + CatBoost 30\% & 0.489822 & 0.498177 & 0.344425 \\
Graph + CatBoost 30\% & 0.491439 & 0.500394 & 0.339809 \\
Graph04D & 0.494881 & 0.503346 & 0.340288 \\
CatBoost d5/lr.06/l2=3 & 0.515282 & 0.523869 & 0.370873 \\
CatBoost d4/lr.03/l2=3 & 0.516796 & 0.525414 & 0.371216 \\
CatBoost d4/lr.06/l2=6 & 0.520583 & 0.528819 & 0.373011 \\
\bottomrule
\end{tabular}
\end{table}

El blend 10\% parece superar Local04D en la sexta cifra decimal:
\[
0.4878453167-0.4878418585
\approx3.46\times10^{-6}.
\]
Pero esta diferencia carece de evidencia práctica. El pooled RMSE es peor:
0.495915 frente a 0.495741. Seleccionar el blend sólo por el mínimo de la media de los mismos
16 splits sería un ejemplo clásico de sobreinterpretación de ruido de validación.

\subsection{Gate leave-one-split-out}

Se restringe el universo de selección a:
\[
\{\text{Local},\ \text{CatBoost},\
\text{Local+CB10},\text{Local+CB20},\text{Local+CB30}\}.
\]
Para cada holdout $h$, el método se selecciona por RMSE medio sobre los otros 15 splits y se
evalúa en $h$. El selector escoge Local puro en 8/16 y Local+CB10 en 8/16. El resultado
concatenado es
\[
\overline{\RMSE}_{LOSO}=0.488884,
\qquad
\RMSE_{\mathrm{pooled}}=0.496838.
\]
Ambos son peores que Local04D fijo. Por tanto, la minúscula diferencia de la tabla completa no
sobrevive selección split-excluded.

\subsection{Diversidad de residuos}

Para pares de métodos se calcula correlación de residuos sobre las 656 pseudo-predicciones
target-matched:
\[
\rho(e^{Local},e^{Graph})=0.9707,
\]
\[
\rho(e^{Local},e^{CB})=0.9423,
\qquad
\rho(e^{Graph},e^{CB})=0.9143.
\]
CatBoost es algo más diverso respecto al grafo, pero sus errores siguen altamente
correlacionados. Su pérdida de accuracy standalone es demasiado grande para que esa diversidad
compense de forma robusta.

\subsection{Stress tests}

Graph+CatBoost puede mejorar algunas métricas target/state-random, pero no conserva la mejor
robustez municipality-grouped. Por ejemplo, Graph+CB30 alcanza 0.4914 target-matched y
0.4627 state-random, frente a 0.4949 y 0.4679 del grafo; sin embargo, grouped empeora de
0.5279 a aproximadamente 0.5508. Un peso 10\% casi preserva grouped, pero sigue siendo peor
que Local en el protocolo primario.

\subsection{Decisión}

\begin{decision}
CatBoost no entra como componente obligatorio de 04F. Se conserva como diagnóstico de
diversidad y sensibilidad. La razón no es que boosting sea débil en general: la evidencia
específica del problema muestra que, después de aprovechar localidad espacial, el anchor
CatBoost no aporta una mejora reproducible.
\end{decision}

\section{04D.1. Refinamiento local y de grafos}
\label{sec:checkpoint04d1}
\subsection{Objetivo del refinamiento}

Checkpoint~04D.1 no vuelve a realizar una búsqueda genérica de modelos. Toma como evidencia
previa que LocalRidge y Graph sobrevivieron 04B, reutiliza exactamente las 34 máscaras
congeladas y explora únicamente grados de libertad relevantes para estas dos familias.

El grid principal contiene 626 métodos predeclarados:
\begin{itemize}
\item 300 variantes de grafo;
\item 324 variantes de Local Ridge;
\item dos anchors fijos.
\end{itemize}
La gran cantidad de variantes no se interpreta como 626 hipótesis independientes; comparten
estructura y muchos hiperparámetros adyacentes. Por eso se estudian superficies de estabilidad
y luego se usa selección leave-one-split-out.

\subsection{Topologías de distancia}

Las topologías de distancia combinan geografía y embedding agronómico:
\[
d^{(\eta)}_{ij}
=
\eta\,\widetilde d^{geo}_{ij}
+
(1-\eta)\,\widetilde d^{agro}_{ij},
\]
con variantes denominadas Geo, GeoAgro90, GeoAgro75, GeoAgro50 y GeoAgro25. En la
nomenclatura, GeoAgro25 corresponde a 25\% geografía y 75\% componente agronómica
normalizada.

Local Ridge se prueba sobre C1 agronomic y C4 all deterministic, con topologías Geo,
GeoAgro75 y GeoAgro50. La búsqueda usa
\[
k\in\{12,16,20,24,30,40\},
\]
\[
\alpha\in\{30,100,300\},
\qquad
p\in\{0.5,1,2\},
\]
donde $p$ es la potencia de distancia en
\[
w_{qi}=\frac{1}{(d_{qi}+10^{-6})^p}.
\]

\subsection{Local Ridge query-specific}

Para cada objetivo $q$ se seleccionan únicamente etiquetas visibles:
\[
\mathcal N_k(q)
=
\operatorname*{arg\,sort}_{i\in\Train^\star}
d_{qi}[1:k].
\]
Entonces
\begin{equation}
\hat{\bm\beta}_q
=
\arg\min_{\bm\beta}
\sum_{i\in\mathcal N_k(q)}
w_{qi}
(y_i-x_i^\top\bm\beta)^2
+
\alpha\|\bm\beta\|_2^2.
\label{eq:cp04d-localridge}
\end{equation}

La regularización es esencial: incluso con $k=24$, C4 tiene una dimensión mucho mayor que el
tamaño del vecindario. Ridge no intenta identificar un vector causal de coeficientes; actúa
como interpolador lineal regularizado en el subespacio local. La estabilidad proviene de
$\alpha$ elevado respecto a un OLS local singular.

\subsection{Búsqueda de grafos}

Los grafos directos varían:
\[
k\in\{4,6,8,10,12\},\qquad
\lambda\in\{0.25,0.5,1,2,4,8\}
\]
sobre cinco topologías. La solución sigue \cref{eq:graph-laplacian}. Valores mayores de
$\lambda$ penalizan más la variación de la función sobre aristas:
\[
\bm f^\top L\bm f
=
\frac{1}{2}\sum_{i,j}w_{ij}(f_i-f_j)^2.
\]
Por tanto, el modelo asume smoothness local del rendimiento en el grafo $X$.

\subsection{Corrección de residuos globales}

Además del grafo directo se probó un modelo de corrección de residuos. Primero se ajusta PLS4
sobre etiquetas visibles y se generan residuos cross-fitted dentro de ese mismo conjunto:
\[
r_i^{CF}
=
y_i-\hat y_i^{PLS,-fold(i)}.
\]
Se ajusta después un campo suave sobre el grafo:
\[
\hat r_q
=
\operatorname{GraphPropagate}(r^{CF}),
\]
y se predice
\[
\hat y_q
=
\hat y_q^{PLS,full}+\hat r_q.
\]
El uso de residuos cross-fitted es crítico. Si se propagaran residuos in-sample de un PLS
ajustado sobre los mismos puntos, su magnitud y estructura estarían contaminadas por fit
optimista.

\subsection{Ranking target-matched}

\begin{table}[H]
\centering
\caption{Mejores configuraciones del refinamiento 04D.1.}
\label{tab:cp04d-ranking}
\small
\begin{tabular}{lrrr}
\toprule
Método & RMSE medio & mediana & peor split \\
\midrule
C4 Geo $k=24,\alpha=30,p=1$ & \textbf{0.487845} & 0.468742 & 0.694700 \\
C1 Geo $k=30,\alpha=30,p=1$ & 0.488302 & 0.4694 & 0.6921 \\
C1 Geo $k=24,\alpha=30,p=1$ & 0.4885 & 0.4684 & 0.6953 \\
C4 Geo $k=30,\alpha=30,p=1$ & 0.4887 & 0.4715 & 0.6925 \\
C4 GeoAgro75 $k=24,\alpha=30,p=1$ & 0.4892 & 0.4718 & 0.6940 \\
\bottomrule
\end{tabular}
\end{table}

La proximidad de varias configuraciones alrededor de $k=20$--30 y $\alpha=30$ es más
importante que el orden exacto de las primeras filas: existe una cuenca estable, no un spike
aislado. Esta observación aumenta la credibilidad del finalista.

\safefigure{../../reports/checkpoint_04d1/figures/local_k_profile.png}
{Perfil de desempeño de Local Ridge al variar el tamaño del vecindario.}
{fig:cp04d-local-k}

\subsection{El mejor modelo robusto no coincide con el mejor target-matched}

El finalista de grafo más robusto es
\artifact{GraphDirect_GeoAgro25_k6_lam8}:
\[
\begin{array}{lr}
\text{target-matched} & 0.494881\\
\text{state-random} & 0.467936\\
\text{state-stratified} & 0.479161\\
\text{municipality-grouped} & 0.527883.
\end{array}
\]

El mejor Local Ridge obtiene 0.487845 target-matched, pero municipality-grouped cerca de
0.687. La diferencia expresa dos objetivos estadísticos:
\begin{itemize}
\item Local Ridge aprovecha intensamente etiquetas geográficamente cercanas, condición que
el matching indica que se parece al test real;
\item GraphDirect conserva mejor comportamiento cuando se elimina un municipio completo.
\end{itemize}

No existe contradicción en conservar ambos como finalistas: uno maximiza la simulación primaria
del test fijo; el otro es un hedge de extrapolación espacial.

\safefigure{../../reports/checkpoint_04d1/figures/refinement_ranking.png}
{Ranking del search local/grafo sobre pseudo-competiciones target-matched.}
{fig:cp04d-refinement}

\subsection{Graph residual versus graph directo}

Las variantes de corrección de residuos PLS quedan alrededor de 0.537 target-matched y
0.679--0.684 grouped, claramente por debajo del grafo directo. Esto sugiere que, en este
problema, la estructura espacial no actúa sólo como una pequeña corrección de un anchor global:
modelar directamente la función sobre el grafo es más competitivo.

\subsection{Leave-one-split-out para selección de método}

Para cada split $h$:
\[
\hat m_{-h}
=
\arg\min_m
\frac{1}{15}
\sum_{s\neq h}\RMSE_{s,m},
\]
y luego se calcula $\RMSE_{h,\hat m_{-h}}$. El resultado es:
\[
\overline{\RMSE}_{LOSO,\mathrm{method}}=0.488366,
\qquad
\RMSE_{\mathrm{pooled}}=0.496137.
\]
Quince de dieciséis selecciones eligieron el C4 Geo $k=24,\alpha=30,p=1$ y una eligió C1
Geo $k=30,\alpha=30,p=1$. Esta concentración confirma que el optimum no depende de un solo
split.

\subsection{Routing por soporte y caveat posterior}

04D.1 también ejecutó selección LOSO por tier y reportó:
\[
\overline{\RMSE}_{LOSO,\mathrm{tier}}=0.485815,
\qquad
\RMSE_{\mathrm{pooled}}=0.494213.
\]
Numéricamente es mejor que la regla única. Sin embargo, una auditoría posterior detectó una
asimetría importante: el universo del selector LOSO de tiers podía incluir los 626 métodos
predeclarados, mientras la propuesta actual para las 59 parcelas se construía sobre un
subconjunto de finalistas. Una ganancia de aproximadamente 0.0026 no justifica ignorar esa
diferencia de candidate universe.

Además, los 16 pseudo-splits reutilizan parcelas; no son 16 datasets independientes. La
selección split-excluded remueve optimismo de mismo split, pero no convierte los resultados
en replicaciones estadísticamente independientes.

\begin{decision}
El routing por soporte se conserva como evidencia exploratoria, no como regla final. La
solución final debe usar un candidate universe explícito y estable o elegir una regla única
suficientemente robusta.
\end{decision}

\subsection{Candidato para las 59 parcelas}

04D.1 materializa predicciones candidatas, incluyendo Local C4 y GraphDirect. Estos valores no
se promueven todavía a entrega: sirven para medir discrepancia entre inductive biases y para
que 04E/04C evalúen si evidencia externa o global puede reducir error antes del congelamiento
de 04F.

\section{04E.1. Evidencia externa SIAP localizada}
\label{sec:checkpoint04e1}
\subsection{Hipótesis externa}

El fuerte Moran's $I$ y la sensibilidad municipal sugieren que una estadística pública
localizada podría aportar un prior regional. SIAP 2025 es especialmente interesante porque es
contemporáneo y público. La hipótesis de 04E.1 no es que SIAP sea equivalente al target
parcelario, sino que
\[
y_i = g(z_{m(i),2025}) + r_i,
\]
donde $z_{m(i),2025}$ es una señal municipal y $r_i$ una corrección parcelaria, podría tener
menor varianza que modelar $y_i$ sin prior.

Para probarla de manera semánticamente correcta se vuelve a la tabla SIAP detallada, evitando
la agregación legada que mezclaba ciclos/modalidades.

\subsection{Scope exacto y fallbacks explícitos}

El scope principal se define como:
\[
\text{Cebada grano}
+\text{Primavera--Verano}
+\text{Temporal}
+\text{CVEGEO}
+\text{2025}.
\]
La elección de Primavera--Verano se alinea con el ciclo abril--octubre del reto. En este
documento se mantiene como decisión operacional del proyecto, no como una validación externa
independiente de nomenclatura SIAP.

La auditoría conserva scopes progresivamente más amplios:
\begin{description}
\item[exact] ciclo y modalidad exactos;
\item[cycle all-mode] ciclo exacto, modalidad libre;
\item[temporal all-cycle] modalidad Temporal, ciclo libre;
\item[all-grain] toda Cebada grano.
\end{description}
Los fallbacks se nombran en el artefacto; nunca se mezclan silenciosamente.

\begin{table}[H]
\centering
\caption{Auditoría de scopes SIAP en 04E.1.}
\label{tab:cp04e-scope}
\begin{tabular}{lrrr}
\toprule
Scope & filas históricas & municipios 2025 & filas 2025 \\
\midrule
exact & 924 & 83 & 83 \\
cycle all-mode & 925 & 83 & 93 \\
temporal all-cycle & 929 & 83 & 83 \\
all-grain & 930 & 83 & 94 \\
\bottomrule
\end{tabular}
\end{table}

La cobertura exacta alcanza 197/197 parcelas. El prior seleccionado tiene cobertura 197/197
antes de fallback por mediana, incluyendo 59/59 objetivos reales. Por tanto, cualquier fracaso
predictivo no puede atribuirse a missingness de SIAP.

\safefigure{../../reports/checkpoint_04e1/figures/siap_scope_coverage.png}
{Cobertura de los scopes SIAP auditados.}
{fig:cp04e-coverage}

\subsection{Construcción del rendimiento municipal}

En un municipio-año, la agregación evita promediar rendimientos reportados sin peso cuando
existen varias filas. Se conservan tanto el rendimiento reconstruido
\[
Y^{prod/harv}_{m,t}
=
\frac{\sum_r \mathrm{VolumenProduccion}_{r}}
{\sum_r \mathrm{Cosechada}_{r}},
\]
como el rendimiento reportado ponderado por superficie cuando corresponde. También se deriva
tasa de siniestro. El objetivo es mantener una definición reproducible aun cuando el source
tenga varias categorías administrativas.

\subsection{Usos evaluados del prior}

Se probaron cuatro formas principales.

\paragraph{Directo.}
\[
\hat y_i=z_{m(i),2025}.
\]

\paragraph{Calibración affine.}
Sobre etiquetas visibles:
\[
y_i=\beta_0+\beta_1z_i+\varepsilon_i,
\]
con Ridge sobre $(1,z_i)$.

\paragraph{Local residual.}
Se define
\[
r_i=y_i-z_i
\]
para etiquetas visibles y se ajusta Local Ridge para $r_q$:
\[
\hat y_q=z_q+\hat r_q.
\]

\paragraph{Graph residual.}
Los residuos visibles se propagan sobre el grafo:
\[
\hat y_q=z_q+\operatorname{Graph}(r)_q.
\]

Finalmente, se prueban blends fijos del 25, 50 y 75\% contra los baselines 04D.

\subsection{Resultados target-matched}

\begin{table}[H]
\centering
\caption{Resultados principales de SIAP 04E.1.}
\label{tab:cp04e-results}
\small
\begin{tabular}{lrr}
\toprule
Método & RMSE medio & pooled MAE \\
\midrule
Local04D & \textbf{0.4878} & 0.3425 \\
Local + SIAP-local 25\% & 0.4922 & 0.3528 \\
Graph04D & 0.4949 & \textbf{0.3403} \\
Graph + SIAP-graph 25\% & 0.5014 & 0.3500 \\
SIAP local residual & 0.5456 & 0.4151 \\
SIAP graph residual & 0.5558 & 0.3976 \\
SIAP affine & $\approx0.837$ & $\approx0.717$ \\
SIAP directo & 1.4368 & 1.3095 \\
\bottomrule
\end{tabular}
\end{table}

La cobertura perfecta no se traduce en resolución parcelaria. En las 138 etiquetas completas:
\[
\RMSE(z,y)=1.4926,
\qquad
\rho_{\mathrm{Spearman}}(z,y)=-0.2181.
\]
La feature SIAP 2025 legada produce RMSE directo 1.4894, prácticamente igual de pobre. Filtrar
correctamente ciclo/modalidad mejora semántica, no garantiza señal parcelaria.

\safefigure{../../reports/checkpoint_04e1/figures/siap_proxy_vs_yield.png}
{Proxy SIAP municipal frente al rendimiento observado de parcela.}
{fig:cp04e-siap-yield}

\subsection{LOSO y evidencia contra la incorporación}

El selector leave-one-target-matched-split-out dispone de Local04D, Graph04D y las variantes
SIAP. En 16/16 holdouts elige Local04D. Esta unanimidad es evidencia más fuerte que comparar
el mejor promedio de una tabla completa: cada holdout se evalúa después de seleccionar con los
otros splits.

\subsection{Resultado grouped que sí favorece SIAP}

El fracaso primario no es uniforme. Bajo municipality-grouped:
\[
\RMSE(\mathrm{Graph04D})=0.5279,
\]
mientras un blend 25\% Graph/SIAP llega a aproximadamente
\[
0.5096.
\]
Esto tiene interpretación coherente: cuando se elimina un municipio completo, un prior
municipal externo puede suplir parte de la información que ya no existe en los labels
cercanos.

Sin embargo, el test real no es un leave-one-municipality-out puro. Las máscaras
target-matched son la simulación diseñada para la geometría específica de los 59 objetivos, y
allí el mismo blend empeora. Por ello el resultado grouped se conserva como stress evidence,
no como autorización para introducir SIAP al predictor final.

\safefigure{../../reports/checkpoint_04e1/figures/external_method_ranking.png}
{Ranking de métodos con evidencia SIAP en el protocolo target-matched.}
{fig:cp04e-ranking}

\subsection{Lección de resolución espacial}

Este checkpoint ilustra un principio general de data fusion: la contemporaneidad no compensa
una discordancia de resolución. SIAP observa un outcome administrativo municipal; el target es
parcelario. Si dentro de un municipio existe heterogeneidad considerable, el prior
$z_m$ tiene un error irreducible para parcelas individuales.

La evidencia externa sigue siendo metodológicamente valiosa porque:
\begin{itemize}
\item verifica una hipótesis razonable;
\item documenta el significado exacto de SIAP;
\item cuantifica su utilidad bajo extrapolación municipal;
\item evita que el equipo incorpore el proxy final por intuición.
\end{itemize}

\begin{decision}
SIAP se conserva como fuente pública documentada y señal de stress. No recibe peso positivo
obligatorio en 04F.
\end{decision}

\section{04F. Reconstrucción final y congelamiento del vector de 59 rendimientos}
\label{sec:checkpoint04f}
\subsection{Principio de cierre}

Checkpoint~04F no debe convertirse en una última búsqueda de hiperparámetros. Llegado este
punto, Local04D, Graph04D, SIAP y CatBoost ya han sido evaluados bajo las mismas
pseudo-competiciones. La función de 04F es congelar una regla reproducible y demostrar que
cualquier alternativa aparentemente mejor supera controles de selección suficientemente
estrictos.

El finalista principal es
\[
\boxed{
\text{\artifact{LocalRidge_C4_all_deterministic_Geo_k24_a30_p1}}
}.
\]
Su nombre codifica:
\begin{itemize}
\item representación C4 all deterministic;
\item vecindad definida por distancia geográfica;
\item $k=24$ etiquetas visibles por consulta;
\item Ridge $\alpha=30$;
\item pesos de distancia inversa con potencia $p=1$.
\end{itemize}

\subsection{Provenance de las 59 predicciones}

Antes de escribir el archivo final, 04F compara las predicciones Local04D materializadas en
04D.1 y 04E.1. La diferencia absoluta máxima es
\[
\max_{i\in\Target}
\left|
\hat y^{04D}_{i,Local}
-
\hat y^{04E}_{i,Local}
\right|
=0.
\]
La misma prueba para Graph04D produce exactamente 0. Esto garantiza que no existe una variante
oculta del modelo bajo el mismo alias y que el vector congelado es trazable a un método
específico.

\subsection{Performance congelada}

Para Local04D en los 16 target-matched splits:
\[
\overline{\RMSE}=0.4878453167,
\qquad
\RMSE_{\mathrm{pooled}}=0.4957414287,
\]
\[
\MAE_{\mathrm{pooled}}=0.3425011445,
\qquad
R^2_{\mathrm{pooled}}=0.6543403.
\]

El peor split tiene RMSE aproximadamente 0.6947. Reportar sólo la media escondería esa
heterogeneidad; por eso el documento conserva mediana, peor split y diagnósticos de soporte.

\subsection{Última hipótesis: blend Local/Graph}

Dado que Graph04D es más robusto bajo municipality-grouped, se evalúa una familia simple:
\[
\hat y^{(w)}
=
(1-w)\hat y^{Local}
+
w\hat y^{Graph}.
\]
Sobre toda la tabla target-matched:
\begin{table}[H]
\centering
\caption{Sensibilidad final Local/Graph.}
\label{tab:cp04f-blends}
\begin{tabular}{rrrrr}
\toprule
$w_{Graph}$ & RMSE medio & pooled RMSE & pooled MAE & peor split \\
\midrule
0.25 & \textbf{0.486768} & \textbf{0.494892} & 0.340479 & 0.693090 \\
0.30 & 0.486780 & 0.494943 & 0.340151 & 0.692886 \\
0.20 & 0.486832 & 0.494914 & 0.340816 & 0.693334 \\
0.40 & 0.487032 & 0.495268 & 0.339582 & 0.692595 \\
0.10 & 0.487188 & 0.495180 & 0.341582 & 0.693939 \\
0.50 & 0.487588 & 0.495887 & \textbf{0.339077} & 0.692461 \\
0.00 & 0.487845 & 0.495741 & 0.342501 & 0.694700 \\
1.00 & 0.494881 & 0.503346 & 0.340288 & 0.694149 \\
\bottomrule
\end{tabular}
\end{table}

La superficie es suave y tiene un mínimo alrededor de 20--30\% de grafo. Si se observara sólo
esta tabla, el 75/25 parecería preferible a Local puro por aproximadamente 0.00108 RMSE medio.

\subsection{Por qué no se usa el 75/25 en el archivo final}

El peso óptimo anterior fue observado usando las mismas 16 pseudo-competiciones que se están
resumiendo. Para cuantificar el optimismo de selección, se ejecuta un leave-one-split-out:
en cada holdout $h$ se selecciona $w$ usando los otros 15 splits y se aplica al holdout.

Los pesos escogidos fueron:
\[
w=0.20\text{ en 3 splits},\quad
w=0.25\text{ en 5},\quad
w=0.30\text{ en 7},\quad
w=0.40\text{ en 1}.
\]
El resultado concatenado:
\[
\overline{\RMSE}_{LOSO}=0.4879409616,
\qquad
\RMSE_{\mathrm{pooled,LOSO}}=0.4959874245.
\]
Esto es ligeramente peor que Local puro. El experimento no prueba que el peso óptimo real sea
cero; demuestra que la evidencia disponible no es suficientemente fuerte para justificar
cambiar una regla ya estable por una mejora de desarrollo de aproximadamente una milésima.

\subsection{Principio de conservadurismo estadístico}

La decisión final puede expresarse como una penalización implícita de complejidad de selección.
Entre dos reglas:
\[
\mathcal R_0=\text{Local fijo},
\qquad
\mathcal R_1=\text{seleccionar peso Local/Graph},
\]
$\mathcal R_1$ tiene un grado de libertad adicional. Para promoverla se exige que su mejora
sobreviva selección fuera del split. No ocurre.

El mismo argumento descarta:
\begin{itemize}
\item el routing por soporte de 04D.1, por la asimetría del candidate universe;
\item Local+CatBoost10, por diferencia de $3.5\times10^{-6}$ que no sobrevive LOSO;
\item SIAP, porque LOSO escoge Local 16/16.
\end{itemize}

\subsection{Discrepancia Local versus Graph como incertidumbre estructural}

Aunque Graph no entra en el valor final, la diferencia
\[
u_i=
|\hat y_i^{Local}-\hat y_i^{Graph}|
\]
es un diagnóstico útil de sensibilidad a inductive bias. En las 59 parcelas:
\[
\bar u=0.1056\ \mathrm{t/ha},
\qquad
\operatorname{median}(u)=0.0769\ \mathrm{t/ha},
\qquad
\max u=0.4980\ \mathrm{t/ha}.
\]
El máximo corresponde a AGC\_003, justamente una parcela que 04A había identificado como
bajo soporte. Esta coherencia es informativa, pero no se usa para cambiar manualmente el
rendimiento: hacerlo después de observar discrepancia sin una regla validada reabriría
overfitting de decisión.

\subsection{Archivo canónico}

El output final es
\begin{center}
\artifact{reports/checkpoint_04f/final_predictions.csv}
\end{center}
con exactamente 59 filas y dos columnas:
\[
\artifact{ID_POLIGONO},\qquad
\artifact{RENDIMIENTO_T_HA}.
\]
No se redondea artificialmente en el artefacto; se conserva precisión completa. El archivo
\artifact{final_prediction_diagnostics.csv} es separado y contiene soporte, Graph04D y
discrepancia. Esta separación impide confundir una variable de diagnóstico con el valor que
se entregará.

\subsection{Estado final de las hipótesis}

\begin{table}[H]
\centering
\caption{Estado de hipótesis al cierre de Checkpoint 04.}
\label{tab:cp04f-hypotheses}
\small
\begin{tabular}{p{0.30\linewidth}p{0.23\linewidth}p{0.36\linewidth}}
\toprule
Hipótesis & Estado & Evidencia \\
\midrule
Localidad geográfica & retenida & mejora target-matched; Moran y pairwise support \\
Grafo semi-supervisado & retenido como diagnóstico & fuerte robustez grouped \\
Temporal-only neighbors & no prioritario & asociación pairwise débil \\
SIAP municipal 2025 & diagnóstico/stress & 16/16 LOSO prefieren Local \\
CatBoost global & diagnóstico & standalone peor; blend no sobrevive gate \\
Routing por soporte & exploratorio & pequeña ganancia con caveat de candidate universe \\
Blend Local/Graph & sensibilidad & mínimo full-table, no mejora LOSO \\
Local C4 Geo k24 & \textbf{final} & mejor regla simple estable \\
\bottomrule
\end{tabular}
\end{table}

\begin{decision}
La predicción final se congela a Local04D puro. A partir de este checkpoint, una nueva
búsqueda sólo tendría sentido si aparece una fuente o hipótesis nueva y explícita; no se
reabre tuning para perseguir diferencias de milésimas sobre los mismos pseudo-splits.
\end{decision}

\chapter{Síntesis técnica, limitaciones y lectura operativa}
\label{chap:synthesis}
\section{Evolución de la evidencia a través de los checkpoints}

La solución final no surge de seleccionar el mínimo de una única tabla. Cada checkpoint reduce
un conjunto distinto de incertidumbres.

Checkpoint~01 reduce incertidumbre \emph{semántica}: qué significa cada fila, qué unidad tiene
cada raster, qué identificador puede usarse para unir fuentes y qué información está realmente
disponible. Checkpoint~02 reduce incertidumbre \emph{representacional}: qué resumen parcelario
se puede construir sin destruir provenance y qué familias muestran shift o sensibilidad.
Checkpoint~03 reduce incertidumbre \emph{de modelo global}: qué representaciones compactas y
qué inductive biases sobreviven nested CV. Checkpoint~03D vuelve después sobre esa misma
pregunta con mayor capacidad y demuestra que el tamaño del grid no era la única limitación.
Checkpoint~04 reduce incertidumbre
\emph{específica del conjunto objetivo}: qué significa la geometría de los 59 $X$ conocidos y
qué métodos transfieren mejor rendimiento hacia esos puntos.

Esta secuencia puede resumirse como:
\[
\text{semántica}
\rightarrow
\text{representación}
\rightarrow
\text{generalización}
\rightarrow
\text{transducción}
\rightarrow
\text{combinación global--local}
\rightarrow
\text{confirmación}
\rightarrow
\text{congelamiento}.
\]

\section{Qué información terminó siendo decisiva}

El hallazgo estadístico dominante es espacial. La asociación de distancia geográfica con
diferencia absoluta de rendimiento, $\rho=0.4882$, supera ampliamente a la distancia C1/C4 y
a la similitud temporal mensual. Moran's $I=0.6797$ confirma que el patrón no es únicamente
una correlación de pares: existe clustering espacial del rendimiento.

Esta estructura explica simultáneamente tres resultados que, vistos por separado, podrían
parecer desconectados:
\begin{enumerate}
\item los modelos globales de 03C tienen buen state RMSE pero grouped RMSE mucho peor;
\item LocalRidge mejora drásticamente en pseudo-competiciones donde existen vecinos
representativos;
\item GraphDirect conserva mucho mejor performance cuando se retienen municipios completos.
\end{enumerate}

El incumbent al cierre de Checkpoint~04 es, por tanto, local no porque los modelos locales
sean universalmente mejores, sino porque la topología observada del reto hace valiosa la
transferencia desde vecinos geográficos. Checkpoint~05 prueba de manera explícita si esa
ventaja local puede complementarse con expertos globales sin perder la disciplina de
validación.

\section{Qué información no fue decisiva}

\subsection{Más variables no fue equivalente a más señal}

B2 empirical completo, B5 all y otras representaciones masivas no dominan consistentemente a
C0/C1. En small-$n$, la dimensionalidad adicional puede aumentar varianza, diluir métricas de
distancia o requerir regularización más fuerte.

\subsection{La temporalidad mensual aislada fue débil}

Las trayectorias satelitales son valiosas para construir features de fenología y estado del
cultivo, pero la regla nearest-neighbor temporal aislada no tiene suficiente asociación con
$|y_i-y_j|$ para desplazar geografía. Esto no equivale a negar valor a remote sensing:
gran parte de C4 y de la capa agronómica procede precisamente de esas series.

\subsection{El proxy municipal contemporáneo no sustituyó una etiqueta parcelaria}

SIAP 2025 tiene cobertura completa y gran relevancia contextual, pero su resolución
administrativa es demasiado gruesa para actuar como target directo. El resultado es una
lección de multi-resolution data fusion: una fuente contemporánea sólo ayuda si su escala
espacial es compatible o si existe una corrección local suficientemente informativa.

\subsection{Compute no fue la restricción final}

Checkpoint~03D prueba esta afirmación directamente. El benchmark balanced ejecutó 230 outer
fits en 248.085 minutos y comparó CatBoost, XGBoost, LightGBM, ExtraTrees,
HistGradientBoosting y controles sobre cuatro representaciones deterministas.

El mejor grouped global pasó de aproximadamente 0.748 en E123/E13 a 0.715339 con HistGB G1.
Sin embargo, ese mismo modelo obtuvo state RMSE 0.549924. XGBoost G1 produjo
0.525847/0.733080. Variantes G2/G3 de CatBoost y ExtraTrees llegaron a state RMSE cercano o
inferior a 0.51, pero se degradaron hasta aproximadamente 0.86--0.93 bajo grouped.

El cómputo adicional mueve una parte de la frontera, pero no elimina el gap entre protocolos
ni crea observaciones independientes nuevas. Con $n=138$, más capacidad no sustituye una
representación defendible y una validación que refleje la estructura espacial.

\section{Lectura de las métricas del incumbent}

El RMSE target-matched de 0.487845 de Local04D debe interpretarse como una estimación interna bajo el
protocolo diseñado para parecerse al conjunto objetivo, no como una garantía del RMSE oculto
de FIRA. Las 16 máscaras comparten parcelas y no constituyen 16 realizaciones independientes
del proceso generador.

Por ello el reporte conserva simultáneamente:
\begin{itemize}
\item RMSE medio por split;
\item RMSE pooled;
\item peor split;
\item protocolos state-random/state/grouped;
\item estabilidad de hiperparámetros;
\item selección leave-one-split-out;
\item discrepancia entre finalistas.
\end{itemize}

La evidencia se considera fuerte cuando varias de estas vistas apuntan en la misma dirección,
no cuando una métrica aislada mejora por una cantidad marginal.

\section{Limitaciones estadísticas}

\subsection{Tamaño muestral}

La limitación fundamental es que sólo existen 138 rendimientos observados. Ningún número de
features o seeds crea observaciones independientes adicionales. Las inferencias sobre
diferencias menores a pocas milésimas deben tratarse con extrema cautela.

\subsection{Dependencia espacial}

Las parcelas no son IID. La propia razón por la que LocalRidge funciona implica que una
interpretación IID de la varianza de CV sería incorrecta. Municipality-grouped funciona como
stress test, pero también representa un dominio más difícil que el conjunto real en algunas
regiones.

\subsection{Dependencia entre pseudo-splits}

Los 16 target-matched repeats reutilizan etiquetas. El leave-one-split-out impide escoger una
regla con el mismo split que luego se califica, pero la información de parcelas se solapa
entre los otros 15 y el holdout. El resultado reduce optimismo de selección, no produce un
estimador completamente independiente.

\subsection{Resolución temporal}

Las métricas temporales de 04A usan trayectorias mensuales canónicas. BASIC/PRO contienen
capturas más finas; una investigación posterior podría usar secuencias de mayor resolución.
No se continuó porque la temporalidad mensual ya mostró una transferencia mucho menor que la
geografía y el objetivo del reto exigía priorizar hipótesis con mejor evidencia.

\subsection{SIAP}

El scope exacto se implementó correctamente a nivel de grano/ciclo/modalidad/CVEGEO, pero la
alineación del ciclo Primavera--Verano con abril--octubre es una decisión del proyecto.
Además, un rendimiento municipal puede incorporar composición de productores, manejo y
superficies distintas de las parcelas del reto.

\section{Limitaciones causales}

Ningún resultado del pipeline debe leerse como estimación causal. Si SOC$\times$CEC,
precipitación o NDVI aparecen en una representación útil, no se concluye que intervenir sobre
esa variable causaría la variación observada de rendimiento. El objetivo es predicción/
reconstrucción bajo datos observacionales.

La autocorrelación espacial también puede representar múltiples mecanismos latentes:
agroclima, manejo, variedad, calendario, estructura edáfica o procesos administrativos. El
grafo es un modelo de dependencia predictiva, no una identificación causal.

\section{Reproducibilidad y separación entre investigación y entrega}

El repositorio conserva explícitamente:
\begin{itemize}
\item configuraciones YAML por checkpoint;
\item código reusable bajo \artifact{src/geocebada/};
\item runners bajo \artifact{tools/};
\item artefactos generados bajo \artifact{reports/};
\item handoffs e historial metodológico;
\item registro de asistencia de IA;
\item un archivo final de predicciones separado de diagnósticos.
\end{itemize}

La entrega no debe construirse copiando valores desde una gráfica ni desde este PDF. Mientras
Checkpoint~05 permanezca pendiente, el artefacto canónico es el CSV de 04F con precisión
completa. Después de ejecutar 05, la fuente canónica será
\artifact{reports/checkpoint_05/final_predictions.csv}; ese archivo contendrá Local04D si
ningún challenger supera todos los gates, o el challenger confirmado si la promoción pasa.

\section{Incumbent al entrar a Checkpoint 05}

La cadena de evidencia de Checkpoints~01--04 produjo el incumbent:
\[
\boxed{
\widehat{\bm y}_{\Target}^{(04F)}
=
\operatorname{LocalRidge}_{C4,\ Geo,\ k=24,\ \alpha=30,\ p=1}
(X_{\Train},y_{\Train},X_{\Target})
}.
\]

La complejidad de los checkpoints justifica por qué esta regla relativamente simple sobrevivió
representaciones de alta dimensión, árboles, CatBoost, PLS, kernels, grafos, priors externos,
routing y blends. No obstante, Checkpoint~05 plantea una hipótesis diferente a las ya cerradas:
no buscar un mejor experto aislado, sino estimar si expertos globales y locales tienen errores
complementarios explotables mediante stacking cross-fitted.

Por tanto, LocalRidge C4 sigue siendo la salida canónica mientras 05 está pendiente. Un
challenger sólo puede reemplazarlo si mejora mean y pooled RMSE en el banco de desarrollo, si
el selector leave-one-development-split-out también mejora ambos criterios y si, después de
quedar fijado, vuelve a superar Local04D en mean y pooled RMSE sobre un banco target-matched
fresco generado $X$-only. El segundo banco no se utiliza para buscar otro challenger.

\section{Siguiente etapa operativa}

Después de resolver el único experimento final de Checkpoint~05, el siguiente bloque del proyecto debe
centrarse en:
\begin{enumerate}
\item empaquetado del archivo de entrega;
\item aplicación/interfaz para explorar parcelas, features y predicciones;
\item explicación de soporte/incertidumbre sin presentar discrepancia como intervalo
probabilístico calibrado;
\item documentación del uso de fuentes externas y de IA;
\item validaciones finales de esquema, IDs y reproducibilidad.
\end{enumerate}

La aplicación podrá mostrar las predicciones de expertos, pesos del mixture-of-experts,
soporte y discrepancia como contexto, pero el valor submission-facing debe leerse del CSV
canónico producido por Checkpoint~05. El bundle
\artifact{models/final/checkpoint05_model.joblib} se reserva para reproducción exacta y para
la futura aplicación Python; una API para parcelas arbitrariamente nuevas constituye un
contrato distinto.


\section{Estado final del modelado}

Checkpoint~05 cerró sin promoción y Checkpoint~03D cerró posteriormente la única duda amplia
de capacidad global. La salida competition-facing permanece Local04D en
\artifact{reports/checkpoint_05/final_predictions.csv}.

Existe un plan 05B, pero es deliberadamente estrecho: sólo puede responder si un peso pequeño
de los nuevos globales 03D añade señal bajo las mismas pseudo-competiciones target-matched.
No es autorización para continuar buscando indefinidamente modelos o hiperparámetros.

El proyecto conserva dos contratos distintos:
\begin{itemize}
\item \textbf{competencia fija-59}: Local04D y el CSV congelado;
\item \textbf{deployment global}: el finalista 03D XGBoost exportado y verificado en ONNX.
\end{itemize}
Confundir ambos contratos produciría una afirmación incorrecta sobre qué modelo generó la
entrega del reto.

\appendix
\chapter{Notación, métricas y formulaciones}
\label{app:notation}
\section{Conjuntos e índices}

\begin{longtable}{p{0.20\linewidth}p{0.68\linewidth}}
\toprule
Símbolo & Definición \\
\midrule
\endhead
$\Train$ & conjunto de 138 parcelas con rendimiento observado. \\
$\Target$ & conjunto de 59 parcelas objetivo con rendimiento oculto. \\
$\Target^\star$ & conjunto pseudo-target ocultado dentro de las 138 etiquetas. \\
$\Train^\star$ & etiquetas visibles de una pseudo-competición,
$\Train\setminus\Target^\star$. \\
$x_i$ & vector de covariables de la parcela $i$. \\
$y_i$ & rendimiento observado de parcela $i$, en toneladas por hectárea. \\
$X_{\Train}$ & matriz de covariables de parcelas etiquetadas. \\
$X_{\Target}$ & matriz observable de covariables de las 59 parcelas objetivo. \\
$W$ o $A$ & matriz de pesos/adyacencia de un grafo. \\
$D$ & matriz diagonal de grados del grafo. \\
$L=D-A$ & Laplaciano no normalizado. \\
$d_{ij}$ & distancia entre parcelas $i,j$ en una topología especificada. \\
$\mathcal N_k(q)$ & $k$ vecinos etiquetados seleccionados para query $q$. \\
\bottomrule
\end{longtable}

\section{Métricas de regresión}

Para observaciones $y_i$ y predicciones $\hat y_i$:
\begin{align}
\MSE
&=
\frac{1}{n}\sum_{i=1}^n(y_i-\hat y_i)^2,\\
\RMSE
&=
\sqrt{\MSE},\\
\MAE
&=
\frac{1}{n}\sum_{i=1}^n|y_i-\hat y_i|,\\
R^2
&=
1-
\frac{\sum_i(y_i-\hat y_i)^2}
{\sum_i(y_i-\bar y)^2}.
\end{align}

En reportes con múltiples splits se distinguen:
\begin{align}
\overline{\RMSE}_{split}
&=
\frac{1}{S}\sum_{s=1}^S\RMSE_s,\\
\RMSE_{\mathrm{pooled}}
&=
\sqrt{
\frac{\sum_s\sum_{i\in s}(y_i-\hat y_i)^2}
{\sum_s n_s}
}.
\end{align}
El primer término da igual peso a cada split; el segundo da igual peso a cada fila
pseudo-predicha. Ambos se conservan porque los tamaños de folds grouped pueden ser
desiguales.

\section{Spearman}

Para pares de variables $u,v$, Spearman es la correlación de Pearson de sus rangos:
\[
\rho_S(u,v)
=
\rho_P(\operatorname{rank}(u),\operatorname{rank}(v)).
\]
En 04A se usa para cuantificar asociación monotónica entre distancia en $X$ y
$|y_i-y_j|$ sin asumir relación lineal.

\section{Moran's \texorpdfstring{$I$}{I}}

Con desviaciones $z_i=y_i-\bar y$ y matriz de pesos $W$:
\[
I
=
\frac{n}{S_0}
\frac{\sum_i\sum_jw_{ij}z_i z_j}
{\sum_i z_i^2},
\qquad
S_0=\sum_i\sum_jw_{ij}.
\]
El $p$ reportado en 04A se obtiene por permutación de valores sobre la topología fija. Moran's
$I$ positivo indica que valores similares tienden a ocupar nodos vecinos
\citep{Moran1950}.

\section{Ridge global y local}

Ridge global:
\[
\widehat\beta
=
\arg\min_\beta
\|y-X\beta\|_2^2+\alpha\|\beta\|_2^2.
\]
Si las columnas están centradas y la matriz es de rango apropiado, la solución cerrada es
\[
\widehat\beta
=
(X^\top X+\alpha I)^{-1}X^\top y.
\]

Local Ridge reemplaza la suma global por vecinos y pesos:
\[
\widehat\beta_q
=
\arg\min_\beta
\sum_{i\in\mathcal N_k(q)}
w_{qi}(y_i-x_i^\top\beta)^2
+\alpha\|\beta\|_2^2.
\]
En el finalista,
\[
k=24,\qquad
\alpha=30,\qquad
w_{qi}=(d^{geo}_{qi}+10^{-6})^{-1}.
\]

\section{PLS}

PLS construye componentes latentes supervisados. A nivel conceptual:
\[
t_h=Xw_h,
\]
con direcciones $w_h$ seleccionadas para capturar covarianza con $y$ sujeto a la deflación de
componentes previos \citep{Wold2001}. Dado que $w_h$ depende de $y$, el fit de PLS debe ocurrir
dentro de train en cada fold.

\section{PCA}

PCA es $X$-only. Para matriz centrada $X$, busca
\[
v_1=\arg\max_{\|v\|=1}v^\top X^\top Xv,
\]
y componentes siguientes ortogonales. En 04A, PCA transductivo puede usar las 197 filas porque
no usa etiquetas y la tarea fija permite conocer $X_{\Target}$.

\section{kNN con pesos}

Para $p\ge0$:
\[
w_{qi}
=
(d_{qi}+\epsilon)^{-p},
\qquad
\hat y_q
=
\frac{\sum_{i\in\mathcal N_k(q)}w_{qi}y_i}
{\sum_{i\in\mathcal N_k(q)}w_{qi}}.
\]
$p=0$ corresponde a vecinos con peso uniforme.

\section{Grafo y Laplaciano}

El grafo kNN usa
\[
A_{ij}
=
\exp\left[-\frac12(d_{ij}/\sigma)^2\right]
\]
para aristas seleccionadas y se simetriza tomando el peso máximo en ambas direcciones.
\[
D=\diag(A\mathbf 1),\qquad
L=D-A.
\]
La energía de suavidad:
\[
f^\top Lf
=
\frac12\sum_{ij}A_{ij}(f_i-f_j)^2.
\]
El estimador de 04B/04D resuelve:
\[
(M+\lambda L+\epsilon I)f
=
M(y-\bar y\mathbf1),
\]
donde $M$ tiene uno en nodos etiquetados visibles y cero en nodos sin etiqueta.

\section{Blend convexo}

Para predictores $a$ y $b$:
\[
\hat y^{(w)}=(1-w)\hat y^{(a)}+w\hat y^{(b)},
\qquad w\in[0,1].
\]
Un mínimo aparente de $w$ sobre la misma validación no se toma automáticamente como regla
final. Checkpoint~04F selecciona $w$ leave-one-split-out antes de decidir si la complejidad
adicional está justificada.

\section{Target-matched pseudo-competition}

El pseudo-target count se deriva de:
\[
n_{\Target^\star}
\approx
|\Train|\frac{|\Target|}{|\Train|+|\Target|}
=
138\frac{59}{197}
\approx41.
\]
La construcción de la máscara sólo usa $X$ y metadata. El score del método se calcula
posteriormente con $y_{\Target^\star}$, ya con la predicción congelada.

\section{Discrepancia de finalistas}

En 04F se conserva
\[
u_i=
|\hat y_i^{Local}-\hat y_i^{Graph}|
\]
como medida de sensibilidad a modelo, no como desviación estándar probabilística. Para
interpretarla como intervalo predictivo sería necesaria calibración específica de cobertura,
que no se realizó.

\chapter{Mapa de artefactos y reproducibilidad}
\label{app:repro}
\section{Arquitectura reproducible}

El proyecto separa configuración, lógica reusable, runners y outputs. Los checkpoints no se
definen por una secuencia informal de notebooks, sino por contratos versionados.

\begin{longtable}{p{0.24\linewidth}p{0.65\linewidth}}
\toprule
Componente & Rol \\
\midrule
\endhead
\artifact{configs/} & contratos YAML, grids, paths, thresholds e invariantes. \\
\artifact{src/geocebada/data/} & loaders y normalización de fuentes. \\
\artifact{src/geocebada/features/} & construcción determinista de capas de features. \\
\artifact{src/geocebada/evaluation/} & split logic, estimadores y métricas de checkpoints. \\
\artifact{src/geocebada/visualization/} & figuras derivadas de artefactos numéricos. \\
\artifact{tools/} & entry points reproducibles para ejecutar cada checkpoint. \\
\artifact{reports/} & métricas, predicciones OOF/pseudo-test, figuras y candidatos. \\
\artifact{docs/} & doctrina, hallazgos canónicos, fuentes e historial. \\
\bottomrule
\end{longtable}

\section{Mapa de ejecución}

La secuencia principal puede reproducirse conceptualmente con:

\begin{lstlisting}[language=bash]
python tools/audit_data_contract_v2.py
python tools/build_parcel_feature_table.py

python tools/build_agronomic_features_v1.py
python tools/run_checkpoint_03b.py
python tools/run_checkpoint_03c.py
python tools/run_checkpoint_03c2.py
python tools/run_checkpoint_03c3.py

python tools/run_checkpoint_04a.py
python tools/run_checkpoint_04b.py
python tools/run_checkpoint_04d1.py
python tools/run_checkpoint_04e1.py
python tools/run_checkpoint_04c1.py
python tools/run_checkpoint_04f.py
python tools/run_checkpoint_05.py

# Additive global-capacity closure, executed after Checkpoint 05:
python tools/run_checkpoint_03d.py
\end{lstlisting}

El orden 04E/04C anterior refleja el orden real de implementación de las hipótesis finales;
la numeración conceptual mantiene 04C como rama global, 04D como local/grafo, 04E como
evidencia externa y 04F como congelamiento.

\section{Artefactos canónicos por checkpoint}

\subsection*{Checkpoint 01}
\begin{itemize}
\item \artifact{configs/data_contract_v2.yaml}
\item \artifact{reports/data_audit_v2.json}
\item \artifact{reports/data_audit_v2.md}
\item \artifact{data/integration_fixture/}
\end{itemize}

\subsection*{Checkpoint 02}
\begin{itemize}
\item \artifact{data/processed/features_v1/parcel_features_all.csv}
\item \artifact{data/processed/features_v1/feature_manifest.json}
\item \artifact{reports/checkpoint_02/cv_folds.csv}
\item \artifact{reports/checkpoint_02/baseline_summary.csv}
\item \artifact{reports/checkpoint_02/checkpoint_02_report.md}
\end{itemize}

\subsection*{Checkpoint 03}
\begin{itemize}
\item \artifact{data/processed/agronomic_features_v1/}
\item \artifact{data/processed/empirical_features_v1/}
\item \artifact{reports/checkpoint_03b/expression_discovery_report.md}
\item \artifact{reports/checkpoint_03c/checkpoint_03c_report.md}
\item \artifact{reports/checkpoint_03c2/checkpoint_03c2_report.md}
\item \artifact{reports/checkpoint_03c3/}
\item \artifact{reports/checkpoint_03d/}
\item \artifact{docs/CHECKPOINT_03D_FINDINGS.md}
\end{itemize}

\subsection*{Checkpoint 04}
\begin{itemize}
\item \artifact{reports/checkpoint_04a/}
\item \artifact{reports/checkpoint_04b/pseudo_competition_splits.csv}
\item \artifact{reports/checkpoint_04d1/}
\item \artifact{reports/checkpoint_04e1/}
\item \artifact{reports/checkpoint_04c1/}
\item \artifact{reports/checkpoint_04f/final_predictions.csv}
\item \artifact{reports/checkpoint_04f/final_prediction_diagnostics.csv}
\end{itemize}

\subsection*{Checkpoint 05}
\begin{itemize}
\item \artifact{reports/checkpoint_05/final_predictions.csv}
\item \artifact{reports/checkpoint_05/final_prediction_diagnostics.csv}
\item \artifact{reports/checkpoint_05/checkpoint_05_report.json}
\item \artifact{docs/CHECKPOINT_05_FINDINGS.md}
\end{itemize}

\subsection*{Checkpoint 03D deployment}
\begin{itemize}
\item \artifact{reports/checkpoint_03d/finalists.csv}
\item \artifact{reports/checkpoint_03d/robustness_summary.csv}
\item \artifact{reports/checkpoint_03d/checkpoint_03d_report.json}
\item local, gitignored: \artifact{models/final/checkpoint03d_global.onnx}
\item local, gitignored: \artifact{models/final/checkpoint03d_global.onnx.json}
\item local, gitignored: \artifact{models/final/checkpoint03d_global.joblib}
\end{itemize}

\section{Invariantes que deben comprobarse antes de una entrega}

\begin{enumerate}
\item El archivo final tiene 59 filas de datos y exactamente dos columnas.
\item \artifact{ID_POLIGONO} es único y coincide exactamente con el set oficial
\artifact{PREDICCION}.
\item No existe ningún target observado en esas 59 filas dentro del input de modelado.
\item La columna final se denomina exactamente \artifact{RENDIMIENTO_T_HA}.
\item Las predicciones Local04D coinciden con los candidatos congelados de 04D.1/04E.1.
\item El CSV final conserva precisión numérica completa.
\item Cualquier versión redondeada del reporte o de la app se trata sólo como presentación.
\end{enumerate}

\section{GPU y determinismo}

CatBoost se ejecuta en GPU cuando corresponde. Las seeds y configuraciones se congelan en YAML.
El predictor final LocalRidge es determinista dados:
\begin{itemize}
\item la misma tabla C4;
\item las mismas coordenadas/IDs;
\item $k=24$, $\alpha=30$, $p=1$;
\item el mismo conjunto de 138 etiquetas.
\end{itemize}

El uso de GPU no forma parte de la semántica del modelo final. Fue una decisión de eficiencia
para los benchmarks CatBoost.

\section{Separación entre artefacto generado y narrativa}

Las tablas de este documento se transcriben de artefactos congelados, pero no constituyen la
fuente de verdad primaria. Si el texto y un CSV/JSON generado difieren, debe investigarse el
runner y corregirse el reporte; nunca debe editarse el output numérico para que coincida con la
narrativa.

\section{Versionado y provenance}

Cada reporte generado registra timestamp y commit cuando el runner lo soporta. La cadena de
provenance final es:
\[
\text{source data}
\rightarrow
\text{Feature Table v1}
\rightarrow
\text{C4 deterministic}
\rightarrow
\text{04B memberships}
\rightarrow
\text{04D LocalRidge}
\rightarrow
\text{04F frozen rule}
\rightarrow
\text{05 promotion gate}
\rightarrow
\text{05 canonical CSV}.
\]

Los Checkpoints~04E y 04C son ramas de validación de hipótesis: pueden generar predicciones
candidatas, pero no alteran silenciosamente la cadena final si sus gates no justifican el
cambio.

\section{Política de outputs pesados}

Los reportes, PDFs y figuras son outputs derivados. En el repositorio, su regeneración pesada
no debe dispararse automáticamente por cada push; los workflows costosos deben ejecutarse de
forma manual y subir artifacts, mientras CI ligero puede revisar tests, lint y estructura.
Esta separación evita convertir una actualización documental en un job de cómputo
innecesario.

\section{Recompilación de este reporte}

Desde \artifact{reporte/tecnico/}:
\begin{lstlisting}[language=bash]
latexmk -pdf -interaction=nonstopmode main.tex
\end{lstlisting}

La bibliografía usa BibLaTeX/Biber. Si no existe latexmk:
\begin{lstlisting}[language=bash]
pdflatex main.tex
biber main
pdflatex main.tex
pdflatex main.tex
\end{lstlisting}

Las figuras usan rutas relativas a \artifact{../../reports/}. La macro
\artifact{\\safefigure} deja un placeholder si una figura no está disponible, permitiendo que
el texto siga compilando en checkouts parciales.


\section{Incidentes reproducibles de los runs finales}

En Checkpoint~05, un nested fold de 18 observaciones hizo inválido
\artifact{PLS(n_components=20)}. La corrección limita componentes al mínimo entre número de
features y tamaño del menor training fold interno. Es una restricción de factibilidad, no
selección basada en target.

En Checkpoint~03D, la primera configuración wide proyectó varios días de cómputo. El run
balanced preservó outer e inner CV y redujo sólo el número de candidatos pesados. Para impedir
mezclar partials incompatibles, \artifact{--resume} verifica un fingerprint del contrato.

El preflight ONNX detectó incompatibilidades reales de conversores. Se fijó
\artifact{onnx<1.22} con \artifact{skl2onnx==1.20.0}; CatBoost usa exportador nativo y
XGBoost/LightGBM negocian el opset máximo soportado. HistGB permanece científicamente válido
aunque no sea deployable con ese stack. El finalista XGBoost exportado pasó el round trip con
diferencia absoluta máxima $3.814697\times10^{-6}$.

\chapter{Catálogo técnico de features y representaciones}
\label{app:features}
\section{Feature Table v1}

La representación base contiene 1,403 predictores trazables a source/family/mode. Para lectura
del equipo, las familias deben entenderse como bloques de evidencia, no como variables
intercambiables.

\subsection{Geometría y administración}

Incluye área oficial, centroides, identificadores y contexto administrativo. Estas columnas son
especialmente importantes para:
\begin{itemize}
\item joins con INEGI/SIAP;
\item cálculo de distancias geográficas;
\item control de folds agrupados;
\item análisis de soporte.
\end{itemize}

Los identificadores administrativos no se interpretan numéricamente. Cuando una variable
categórica se materializa como código, su valor entero no representa una distancia ordinal.

\subsection{BASIC histórico y 2025}

Los resúmenes BASIC se mantienen separados por sensor, índice, estadístico y ventana temporal.
Una feature típica codifica:
\[
\text{source}\times\text{period}\times\text{sensor}
\times\text{index}\times\text{statistic}\times\text{month}.
\]
Esta granularidad explica por qué la tabla base alcanza cientos de columnas aunque sólo existan
197 parcelas.

Las variables históricas 2022--2024 contribuyen al modo clean; los resúmenes 2025 completos
son competition. La distinción permite comparar:
\[
\text{estado estructural/histórico}
\quad\text{vs}\quad
\text{condición del ciclo evaluado}.
\]

\subsection{PRO Planet}

Las features Planet son competition por construcción, al corresponder a 2025. NDVI, EVI, LAI
y MSAVI se resumen de forma consistente con la entidad parcela. No se fusionan con Sentinel-2
para crear un único NDVI universal; las comparaciones inter-sensor aparecen como features
derivadas en 03A/03B.

\subsection{Clima oficial}

Se conservan precipitación, temperatura mínima y máxima históricas y 2025. Las capas
agronómicas derivan:
\[
T_{\mathrm{mean}},\quad DTR,\quad GDD_{\mathrm{proxy}},
\quad P_{\mathrm{early}},P_{\mathrm{mid}},P_{\mathrm{late}}.
\]
La existencia de la fuente base y de una transformación derivada no constituye duplicación
accidental: permite que modelos no lineales aprendan desde inputs crudos y que modelos
lineales dispongan de una base más útil.

\subsection{WaPOR}

Las 22 features de competition representan resúmenes de AETI/NPP y provenance asociado. Los
totales respetan duración de dekads. Interacciones hídricas posteriores se mantienen separadas
de estos valores source-derived.

\subsection{SoilGrids}

Las 90 features SoilGrids de la tabla clean resumen nueve propiedades a distintas
profundidades y estadísticos. Las transformaciones de 03A usan perfiles 0--30 y 30--60 para
crear gradientes y razones, evitando un polynomial expansion indiscriminado.

\subsection{SIAP}

La capa base contiene historia municipal y variables contemporáneas. Deben distinguirse:
\begin{enumerate}
\item SIAP histórico como contexto regional;
\item SIAP 2025 legado de Feature Table v1;
\item SIAP 2025 exact-scope re-auditado en 04E.1.
\end{enumerate}
La tercera es semánticamente superior para auditar la fuente, pero su RMSE directo demuestra
que mayor pureza semántica no implica mayor resolución predictiva.

\section{Agronomic Features v1}

\begin{longtable}{p{0.28\linewidth}rp{0.50\linewidth}}
\toprule
Familia & $p$ & Interpretación \\
\midrule
\endhead
phenology & 162 & forma/timing del ciclo por sensor e índice.\\
phenology anomaly & 41 & contraste 2025 contra baseline histórico parcelario.\\
water productivity & 40 & relaciones entre NPP, AETI, precipitación y crop state.\\
thermal & 36 & temperatura media, DTR y GDD proxies.\\
cross-domain & 19 & interacciones seleccionadas entre clima, suelo, SIAP y satélite.\\
soil profile & 18 & gradientes y razones verticales.\\
water timing & 10 & ventanas de precipitación y distribución estacional.\\
sensor agreement & 9 & concordancia S2--Planet en NDVI/EVI/LAI.\\
soil interaction & 8 & SOC, CEC, textura, N, pH y bases compactas.\\
nonlinear basis & 7 & log-transformaciones y otras bases target-free.\\
\bottomrule
\end{longtable}

Toda feature de esta capa puede escribirse como
\[
z_{ij}=g_j(X_i),
\]
sin dependencia de $y$. Esa propiedad permite materializarla una vez para las 197 parcelas.

\section{Empirical Features v1}

\begin{longtable}{p{0.31\linewidth}rp{0.47\linewidth}}
\toprule
Familia & $p$ & Interpretación \\
\midrule
\endhead
temporal shape & 162 & TV, roughness, curvature, autocorrelación, timing y entropía.\\
aggregate geometry & 72 & resúmenes matemáticos de bloques y distribuciones.\\
symmetric change & 42 & cambio 2025/histórico robusto a denominadores cercanos a cero.\\
short-baseline condition & 41 & posición de 2025 en rango histórico corto.\\
sensor shape similarity & 18 & Pearson/Spearman/cosine/RMSE normalizado inter-sensor.\\
\bottomrule
\end{longtable}

La capa materializada es target-free. El miner supervisado es un componente separado:
\[
\mathcal E_f=A(X_{\mathrm{train}(f)},y_{\mathrm{train}(f)}).
\]
Nunca se escribe un único set $\mathcal E$ descubierto con las 138 etiquetas completas y luego
se presenta como input independiente.

\section{Representaciones modeladas}

\subsection{Checkpoint 03C.1}

Las B-representations sirven para medir el efecto de capas:
\[
B0=\text{base},\quad
B1=\text{agro},\quad
B2=\text{empirical},
\]
\[
B3=\text{base+agro},\quad
B4=\text{base+empirical},\quad
B5=\text{all},\quad
B7=\text{all+fold-local discovery}.
\]

\subsection{Checkpoint 03C.2}

El benchmark anidado reduce el espacio:
\[
C0=\text{base},\quad
C1=\text{agro},
\]
\[
C2=\text{all deterministic + discovery},\quad
C3=\text{base+agro+selected empirical support+discovery}.
\]

\subsection{Checkpoint 04}

La topología transductiva retiene:
\[
C0,\qquad C1,\qquad C4=\text{base+agro+empirical deterministic}.
\]
C4 excluye fórmulas descubiertas con $y$ para que PCA/distancias/grafos puedan aprenderse sobre
las 197 parcelas sin contaminar la construcción $X$-only.

\section{Principio general de provenance}

Una feature derivada debe declarar:
\[
(\text{source},\text{parents},\text{formula},\text{mode},\text{evidence}).
\]
Esta tupla es más importante que el nombre humano de la columna. Dos features con nombres
parecidos pero distinto sensor, periodo o scope administrativo no deben considerarse
equivalentes.


\chapter{Protocolos de validación y selección}
\label{app:validation}
\section{Matriz conceptual de protocolos}

No existe un único score que responda todas las preguntas del proyecto. La
\cref{tab:validation-matrix} resume los protocolos principales.

\begin{longtable}{p{0.22\linewidth}p{0.18\linewidth}p{0.22\linewidth}p{0.27\linewidth}}
\caption{Protocolos de validación y su interpretación.}
\label{tab:validation-matrix}\\
\toprule
Protocolo & Tamaño & X objetivo visible & Pregunta principal \\
\midrule
\endfirsthead
\toprule
Protocolo & Tamaño & X objetivo visible & Pregunta principal \\
\midrule
\endhead
state-stratified & 5 folds sobre 138 & no como target fijo en 02--03 & interpolación con estados representados.\\
municipality-grouped & 5 folds sobre 138 & no como target fijo en 02--03 & extrapolación a municipios retenidos.\\
target-matched & 16 $\times$ 41 & sí, bajo contrato transductivo & aproximación primaria a la geometría de las 59 parcelas.\\
state-random & 8 $\times$ 41 & sí & control aleatorio respetando estado.\\
LOSO de pseudo-splits & 16 holdouts & sí & controlar selección de método/peso contra mismo split.\\
\bottomrule
\end{longtable}

\section{Nested CV de Checkpoint 03}

La nested CV separa dos capas:
\[
\text{inner CV}\rightarrow \text{selección de hiperparámetro},
\]
\[
\text{outer holdout}\rightarrow \text{estimación de desempeño}.
\]
Formalmente, para outer fold $f$:
\[
\hat\lambda_f=
\arg\min_{\lambda\in\Lambda}
\widehat R_{\mathrm{inner},f}(\lambda),
\]
\[
\widehat R_{\mathrm{outer},f}
=
R\left(
y_{V_f},
f_{\hat\lambda_f,T_f}(X_{V_f})
\right).
\]
La primera operación no puede observar $y_{V_f}$. Esto incluye selección de expresiones
supervisadas, número de componentes PCA/PLS cuando se tunearon y parámetros de modelos.

\section{Pseudo-competition de Checkpoint 04}

En una pseudo-competición $s$, definimos:
\[
\Train_s=\Train\setminus\Target_s^\star,
\qquad
|\Target_s^\star|=41.
\]
El modelo puede usar
\[
X_{\Train},X_{\Target}
\]
para transformaciones $X$-only globales, y puede usar
\[
X_{\Target_s^\star}
\]
durante el diseño de topología. Sin embargo, el operador supervisado satisface:
\[
\hat f_s
=
A(X_{\Train_s},y_{\Train_s};X_{\Target_s^\star}),
\]
con
\[
y_{\Target_s^\star}
\]
reservado hasta evaluación.

Esta distinción es la razón por la que un PCA transductivo puede ser correcto mientras una
selección de variables por correlación con todas las 138 etiquetas no lo sería.

\section{Matching de pseudo-targets}

El target matching utiliza únicamente descriptores de $X$ y metadata. Sea $q(S)$ una función
de calidad de una máscara candidata $S$ que combina distancia entre perfiles, composición
municipal y penalización de diversidad. La máscara seleccionada es aproximadamente:
\[
S^\star
=
\arg\min_{S:|S|=41}
q(S;X_{\Target}),
\]
entre un conjunto finito de candidatos generados con seed congelada.

La validez del procedimiento depende de no introducir ninguna estadística derivada de
$y_S$ en $q$. Las diferencias de yield de vecinos observadas en 04A pueden usarse para
diagnóstico científico, pero no para escoger qué etiquetas se ocultan.

\section{Leave-one-pseudo-split-out}

Para candidate set $\mathcal M$, la selección de método en holdout $h$ es:
\[
m_h^\star
=
\arg\min_{m\in\mathcal M}
\frac{1}{S-1}
\sum_{s\neq h}\RMSE_{s,m}.
\]
Entonces:
\[
R_h
=
\RMSE_{h,m_h^\star}.
\]
El resumen LOSO es
\[
\bar R_{\mathrm{LOSO}}=\frac1S\sum_h R_h.
\]

Este procedimiento controla el error más directo de escoger el mejor método usando el mismo
split que se puntúa. No elimina toda dependencia porque las máscaras reutilizan parcelas. Su
interpretación correcta es \emph{split-excluded selection}, no nested CV con $S$ datasets
independientes.

\section{Candidate universe}

Una condición sutil es que $\mathcal M$ debe ser el mismo universo de métodos que la regla
final pretende utilizar. Si el LOSO puede escoger entre 626 métodos pero la regla real sólo
puede escoger entre diez finalistas, la aparente performance del selector y la regla aplicada
no son el mismo procedimiento.

Este caveat motivó no promover el routing LOSO de 04D.1 pese a su RMSE medio 0.4858. En 04C.1
y 04F el candidate universe se redujo explícitamente antes de seleccionar:
\begin{itemize}
\item 04C.1: Local, CatBoost estable y tres blends Local/CatBoost;
\item 04F: ocho pesos Local/Graph predeclarados.
\end{itemize}

\section{Lectura de diferencias pequeñas}

Para dos métodos $a,b$:
\[
\Delta=\overline{\RMSE}_a-\overline{\RMSE}_b.
\]
Cuando $|\Delta|$ es del orden de $10^{-3}$ o menor, el proyecto exige evidencia adicional:
\begin{enumerate}
\item estabilidad entre splits;
\item pooled RMSE compatible;
\item selección LOSO que no revierta la ventaja;
\item ausencia de un incremento desproporcionado de complejidad;
\item comportamiento razonable en stress protocols.
\end{enumerate}

Este criterio evita convertir ruido de selección en una narrativa de mejora. El caso extremo es
Local+CatBoost10: $\Delta$ en mean RMSE es de apenas $3.5\times10^{-6}$ y no se sostiene en
pooled ni LOSO.

\section{Protocolos de estrés versus objetivo primario}

Municipality-grouped no se optimiza como si fuera el score oculto real; se usa para detectar
fragilidad. Una regla puede ser aceptable si:
\[
R_{\mathrm{target}} \text{ es muy bajo}
\]
y
\[
R_{\mathrm{grouped}} \text{ no introduce un riesgo incompatible},
\]
aun si no es el mínimo grouped.

Graph04D se conserva precisamente por esta razón: no gana el objetivo primario, pero ofrece
una referencia robusta que permite detectar parcelas donde Local puede ser especialmente
sensible a la disponibilidad de vecinos.


\chapter{Resultados negativos e hipótesis cerradas}
\label{app:negative}
\section{Por qué documentar resultados negativos}

Con 138 etiquetas, una parte importante del valor científico del proyecto consiste en cerrar
hipótesis. Si sólo se conservaran los experimentos ganadores, futuros integrantes podrían
repetir búsquedas ya realizadas o interpretar la ausencia de una variable como una omisión
accidental. Este apéndice registra hipótesis plausibles que fueron evaluadas y no se
promovieron.

\section{Más dimensionalidad determinista}

La combinación de base, agronomic y empirical alcanza más de 2,000 variables en 03C.1. No
domina consistentemente a representaciones más compactas. El resultado evita la heurística
\[
p_{\mathrm{nuevo}}>p_{\mathrm{viejo}}
\Longrightarrow
\text{modelo mejor}.
\]
En high-$p$/small-$n$, agregar ruido o variables redundantes puede deteriorar regularización,
distancias y estabilidad.

\section{Empirical-only}

B2 empirical es débil como representación aislada en 03C.1: por ejemplo, competition +
ExtraTrees produce OOF RMSE alrededor de 0.647 frente a 0.500 de B0. La conclusión no es que
sus variables sean inútiles. Algunas ayudan como support primitives y C4, pero la capa no
sustituye a la información source/agronómica.

\section{Kernel Ridge no lineal}

PCA+RBF y PCA+polynomial-2 quedaron muy por detrás de PLS/Ridge/trees en nested CV. El
proyecto no siguió expandiendo grado, gamma o dimensiones PCA sólo porque el hardware lo
permitía. La evidencia disponible indicaba que la familia no justificaba el aumento del
search space.

\section{Temporal nearest neighbors como predictor principal}

Pearson, best-lag correlation y DTW de trayectorias mensuales muestran asociación muy débil con
diferencia absoluta de yield. La temporalidad se retiene como feature/contexto, pero no como
la métrica principal de vecinos.

\section{Domain adaptation global agresiva}

El clasificador adversarial train--target produce AUC 0.4892. No existe evidencia de un
separation boundary global fuerte. Por ello, importance weighting global no se convirtió en
la prioridad de 04D.

\section{Graph residual sobre PLS}

La hipótesis
\[
\hat y=\hat y^{PLS}+\hat r^{Graph}
\]
parecía natural a partir de los residuos espacialmente autocorrelacionados bajo grouped CV.
Sin embargo, las variantes de residual graph quedan aproximadamente en 0.537 target-matched,
muy por detrás de GraphDirect y LocalRidge. La estructura espacial útil no parece limitarse a
una corrección residual pequeña de PLS.

\section{Routing por soporte}

Un primer routing de 04B y el LOSO de 04D sugieren que diferentes tiers pueden favorecer
métodos distintos. Sin embargo:
\begin{itemize}
\item el candidate universe de la validación LOSO no coincidía perfectamente con la regla
operacional final;
\item la ganancia era pequeña;
\item las pseudo-competiciones son dependientes.
\end{itemize}
Por ello no se aplican 59 decisiones de modelo específicas por tier en el output final.

\section{SIAP directo y calibrado}

SIAP exact-scope tiene cobertura completa, pero:
\[
\RMSE_{\mathrm{direct,target}}\approx1.4368,
\]
y sobre las 138 etiquetas completas:
\[
\RMSE\approx1.4926,\qquad
\rho_S\approx-0.2181.
\]
La calibración affine mejora respecto al uso directo pero sigue alrededor de 0.837 RMSE. Los
modelos residuales tampoco alcanzan Local/Graph.

La señal sí mejora extrapolación municipality-grouped en un blend de grafo, por lo que no se
declara inútil. La conclusión precisa es: \emph{no aporta valor suficiente al protocolo
primario para recibir peso en la reconstrucción final}.

\section{CatBoost como final anchor}

CatBoost C1 fue competitivo en 03C.2, pero una vez que se introducen métodos locales:
\[
\RMSE_{CB}\approx0.515\text{--}0.521
>
0.487845=\RMSE_{Local}.
\]
Un peso 10\% en Local produce una diferencia en mean RMSE de sólo
$3.5\times10^{-6}$ y pooled RMSE peor. El gate LOSO queda en 0.488884. No hay base para
promoverlo.

\section{Blend Local/Graph seleccionado post hoc}

El 75/25 produce 0.486768 en el conjunto completo de 16 pseudo-splits, nominalmente mejor que
0.487845. Sin embargo, seleccionar el peso usando 15 splits y aplicarlo al dieciseisavo
produce 0.487941. El mínimo full-table se registra como sensibilidad, no como regla final.

\section{Heat excess mensual a \texorpdfstring{$25^\circ$C}{25 C}}

La base
\[
H_{25}
=
\sum_m d_m\max(T_{\mathrm{mean},m}-25,0)
\]
resultó constante en la tabla canónica. Las temperaturas medias mensuales no tienen suficiente
resolución para capturar episodios de calor corto. La feature y su interacción hídrica se
eliminan, mostrando por qué una hipótesis mecanicista debe pasar checks de variación antes de
modelado.

\section{CHIRPS en la capa agronómica v1}

Aunque CHIRPS ofrece precipitación diaria, sólo una feature sobrevivió Feature Table v1 y se
mantuvo una alerta de QC. 03A y varias representaciones 04 excluyen CHIRPS explícitamente.
Esto evita que una fuente de alta resolución aparente añada sofisticación cuando la cadena de
extracción todavía no está suficientemente certificada.

\section{ERA5-Land}

La fuente fue investigada, pero los problemas de descarga/costo del backend impidieron
certificar un artefacto completo. No se mezclan archivos parciales con fuentes completas. Una
fuente no disponible de manera reproducible se trata como hipótesis futura, no como input
oculto del modelo final.

\section{Conclusión de los resultados negativos}

La estrategia no fue probar métodos hasta que alguno mejorara unas milésimas. El proyecto
avanzó cerrando ramas:
\[
\text{hipótesis}
\rightarrow
\text{protocolo}
\rightarrow
\text{resultado}
\rightarrow
\text{decisión}.
\]
Esa trazabilidad es especialmente importante para trabajo en equipo: permite distinguir una
idea aún no probada de una idea ya evaluada y descartada bajo un contrato explícito.


\chapter{Predicciones finales para las 59 parcelas}
\label{app:predictions}
\section{Advertencia de uso}

La tabla de este apéndice reproduce el vector \emph{incumbente} de 59 predicciones
congelado por Checkpoint~04F, redondeado a seis decimales únicamente para lectura humana.
Checkpoint~05 está implementado pero todavía pendiente de ejecución; por tanto, mientras su
gate de promoción no termine, el artefacto canónico sigue siendo
\begin{center}
\artifact{reports/checkpoint_04f/final_predictions.csv}
\end{center}
Cuando Checkpoint~05 termine, este apéndice debe regenerarse desde
\artifact{reports/checkpoint_05/final_predictions.csv} antes de considerar cerrado el reporte,
incluso si el gate retiene Local04D. De ese modo la tabla documentada y el manifest final
pertenecen al mismo checkpoint de cierre.
No se deben copiar manualmente los valores de esta tabla para construir la entrega.

\section{Vector incumbente de Checkpoint 04F}

\begin{longtable}{lr}
\caption{Predicciones incumbentes congeladas por Checkpoint 04F.}
\label{tab:final-predictions}\\
\toprule
\artifact{ID_POLIGONO} & \artifact{RENDIMIENTO_T_HA} \\
\midrule
\endfirsthead
\toprule
\artifact{ID_POLIGONO} & \artifact{RENDIMIENTO_T_HA} \\
\midrule
\endhead
AGC\_003 & 3.619219 \\
AGC\_011 & 4.749113 \\
AGC\_016 & 4.495977 \\
AGC\_025 & 2.965280 \\
AGC\_027 & 2.919605 \\
AGC\_028 & 4.759666 \\
AGC\_029 & 3.117419 \\
AGC\_030 & 3.141298 \\
AGC\_031 & 4.152739 \\
AGC\_036 & 4.381511 \\
AGC\_039 & 2.878491 \\
AGC\_042 & 3.054654 \\
AGC\_043 & 4.385614 \\
AGC\_050 & 2.957521 \\
AGC\_052 & 4.513075 \\
AGC\_061 & 2.930974 \\
AGC\_063 & 4.277952 \\
AGC\_065 & 4.457068 \\
AGC\_071 & 3.088721 \\
AGC\_075 & 3.029695 \\
AGC\_076 & 4.442815 \\
AGC\_082 & 4.449122 \\
AGC\_085 & 4.028177 \\
AGC\_086 & 2.970075 \\
AGC\_087 & 2.798619 \\
AGC\_088 & 4.367917 \\
AGC\_091 & 4.523682 \\
AGC\_093 & 3.021425 \\
AGC\_094 & 4.772291 \\
AGC\_095 & 4.611794 \\
AGC\_107 & 4.496003 \\
AGC\_113 & 2.982389 \\
AGC\_116 & 4.390382 \\
AGC\_118 & 3.262602 \\
AGC\_120 & 2.965691 \\
AGC\_126 & 4.579298 \\
AGC\_128 & 4.481228 \\
AGC\_129 & 3.477291 \\
AGC\_130 & 4.559845 \\
AGC\_136 & 4.513440 \\
AGC\_142 & 4.556583 \\
AGC\_143 & 4.531609 \\
AGC\_147 & 2.962276 \\
AGC\_149 & 4.525898 \\
AGC\_153 & 4.144400 \\
AGC\_154 & 4.565747 \\
AGC\_156 & 3.296195 \\
AGC\_158 & 4.162532 \\
AGC\_173 & 2.994752 \\
AGC\_174 & 4.490803 \\
AGC\_175 & 4.511380 \\
AGC\_177 & 4.268497 \\
AGC\_182 & 4.460818 \\
AGC\_183 & 2.947021 \\
AGC\_186 & 4.244017 \\
AGC\_188 & 4.527555 \\
AGC\_190 & 4.525003 \\
AGC\_191 & 4.801678 \\
AGC\_192 & 4.628534 \\
\bottomrule
\end{longtable}

\section{Resumen descriptivo del vector incumbente}

Estas estadísticas describen las predicciones; no son una evaluación porque los 59 targets
reales permanecen ocultos:

\[
\min \hat y=2.798619,\qquad
\max \hat y=4.801678,
\]
\[
\overline{\hat y}=3.944322,\qquad
\operatorname{median}(\hat y)=4.367917,\qquad
s_{\hat y}=0.709443.
\]

La distribución predicha no se postprocesa para forzar momentos, cuantiles o límites
administrativos. Cualquier clipping posterior sería una nueva regla de modelado y requeriría
evidencia de validación propia.

\section{Relación con el diagnóstico Local--Graph}

El archivo separado \artifact{final_prediction_diagnostics.csv} contiene para cada ID la
predicción final Local04D, la predicción Graph04D, el tier de soporte y su discrepancia
absoluta. Esos campos se excluyen de la tabla submission-facing para impedir cambios manuales
por parcela. La discrepancia media es 0.1056 t/ha, la mediana 0.0769 t/ha y el máximo 0.4980
t/ha.

El máximo se observa en AGC\_003. El hecho de que esta parcela también aparezca entre los
objetivos de bajo soporte es coherente con mayor sensibilidad a inductive bias, pero no
constituye una etiqueta indirecta ni una justificación para alterar su predicción congelada.


\clearpage
\printbibliography[heading=bibintoc,title={Referencias}]
\end{document}
